% programmes-cpge-mp.tex — programme de physique-chimie de la classe de MP
% Texte officiel de 2021 (voie Mathématiques et physique), publié par le ministère de
% l’Enseignement supérieur. Fichier de données du paquet programmes-cpge.
% Format : voir la documentation, « Ajouter un programme ».
% Le texte officiel écrit parfois l’apostrophe avec un guillemet ouvrant (« l‘électrochimie ») :
% coquille corrigée en ’ partout (voir la documentation, « Écarts au texte officiel »).

\cpgefiliere{MP}{
  discipline = {Physique-chimie},
  classe = {MP},
  voie = {Mathématiques et physique},
  annee = {2021},
}

\cpgerubrique{MP/P}{
  parent = {MP},
  numero = {P},
  niveau = {1},
  titre = {Préambule},
  chapeau = {},
}

\cpgerubrique{MP/P/1}{
  parent = {MP/P},
  numero = {P.1},
  niveau = {2},
  titre = {Objectifs de formation},
  chapeau = {
    \cpgeitem{Le programme de physique-chimie de la classe de MP est conçu comme un socle cohérent et ambitieux de connaissances et de capacités scientifiques s’appuyant sur celles déjà travaillées au lycée et en classe de MPSI. Le programme vise à préparer les étudiants à un cursus d’ingénieur, de chercheur ou d’enseignant. Il s’agit de renforcer chez l’étudiant les compétences inhérentes à la pratique de la démarche scientifique : observer et s’approprier, analyser et modéliser, réaliser et valider, et enfin communiquer et valoriser ses résultats.}
    \cpgeitem{L’acquisition de ce socle par les étudiants constitue un objectif prioritaire pour l’enseignant. Parce que la physique et la chimie sont avant tout des sciences expérimentales qui développent la curiosité, la créativité et l’analyse critique, l’expérience est au cœur de son enseignement, que ce soit en cours ou lors des séances de travaux pratiques. Les activités expérimentales habituent les étudiants à se confronter au réel, comme ces derniers auront à le faire dans l’exercice de leur métier.}
    \cpgeitem{De même, l’introduction de capacités numériques dans le programme prend en compte la place nouvelle des sciences numériques dans la formation des scientifiques notamment dans le domaine de la simulation. Elles offrent aux étudiants la possibilité d’effectuer une modélisation avancée du monde réel, par exemple par la prise en compte d’effets non linéaires.}
    \cpgeitem{La démarche de modélisation occupe également une place centrale dans le programme pour former les étudiants à établir, de manière autonome, un lien fait d’allers-retours entre le « monde » des objets, des expériences, des faits, et celui des modèles et des théories. L’enseignant doit rechercher un point d’équilibre entre des approches complémentaires : conceptuelle et expérimentale, abstraite et concrète, théorique et appliquée, inductive et déductive, qualitative et quantitative. La construction d’un modèle passe aussi par l’utilisation maîtrisée des mathématiques dont un des fondateurs de la physique expérimentale, Galilée, énonçait déjà qu’elles sont le langage dans lequel est écrit le monde.}
    \cpgeitem{Enfin, l’autonomie et la prise d’initiative sont spécifiquement développées à travers la pratique d’activités du type « résolution de problèmes » qui visent à exercer les étudiants à mobiliser de façon complémentaire connaissances et capacités pour répondre à un questionnement ou atteindre un but sans qu’aucune démarche de résolution ne soit fournie.}},
}

\cpgerubrique{MP/P/2}{
  parent = {MP/P},
  numero = {P.2},
  niveau = {2},
  titre = {Organisation du programme},
  chapeau = {
    \cpgeitem{Le programme est organisé en deux parties.}
    \cpgeitem{Dans la première partie, intitulée « Formation expérimentale », sont décrits les objectifs de formation sur le thème « Mesures et incertitudes » ainsi que les méthodes et les capacités expérimentales que les étudiants doivent maîtriser à la fin de l’année scolaire. Leur mise en œuvre doit notamment s’appuyer sur des problématiques concrètes identifiées en gras dans la seconde partie du programme intitulée « Contenus thématiques ». Elles doivent être programmées par l’enseignant de façon à assurer un apprentissage progressif de l’ensemble des capacités attendues.}
    \cpgeitem{La seconde partie, intitulée « Contenus thématiques » est structurée autour de neuf thèmes : « Mécanique », « Éléments de traitement du signal », « Optique », « Électromagnétisme », « Thermodynamique », « Physique quantique », « Éléments de thermodynamique statistique », « Transformations chimiques de la matière : aspects thermodynamiques » et « aspects   thermodynamique et cinétique de l’électrochimie ». La présentation en deux colonnes (« notions et contenus » et « capacités exigibles ») met en valeur les éléments clefs constituant le socle de connaissances et de capacités dont l’assimilation par tous les étudiants est requise. Certains items de cette seconde partie, identifiés en caractères gras dans la colonne « capacités exigibles », se prêtent particulièrement à une approche expérimentale. Ils doivent être abordés en priorité lors de séances de travaux pratiques où l’autonomie et l’initiative individuelle de l’étudiant doivent être privilégiées. La présence de capacités numériques explicitées atteste par ailleurs de la volonté de renforcer ce volet de la formation des étudiants.}
    \cpgeitem{Trois annexes sont consacrées d’une part au matériel nécessaire à la mise en œuvre des programmes, d’autre part aux outils mathématiques et aux outils numériques que les étudiants doivent savoir mobiliser de façon autonome dans le cadre des enseignements de physique-chimie à la fin de l’année de la classe de MP.}
    \cpgeitem{Ce programme précise les objectifs de formation à atteindre pour tous les étudiants. Il n’impose en aucun cas une progression ; celle-ci relève de la liberté pédagogique de l’enseignant.}},
}

\cpgerubrique{MP/P/3}{
  parent = {MP/P},
  numero = {P.3},
  niveau = {2},
  titre = {Les compétences travaillées dans le cadre de la démarche scientifique},
  chapeau = {
    \cpgeitem{L’ensemble des activités proposées en classe préparatoire aux grandes écoles – activités expérimentales, résolutions de problèmes, TIPE, etc. – permet de travailler les compétences de la démarche scientifique qui figurent dans le tableau ci-dessous. Des capacités associées sont explicitées afin de préciser les contours de chaque compétence, elles ne constituent donc pas une liste exhaustive et peuvent parfois relever de plusieurs domaines de compétences. L’ordre de présentation de ces compétences ne préjuge pas d’un ordre de mobilisation de ces dernières lors d’une activité.}
    \cpgeitem{Les compétences doivent être acquises à l’issue de la formation en CPGE. Elles nécessitent d’être régulièrement mobilisées par les étudiants et sont évaluées en s’appuyant, par exemple, sur l’utilisation de grilles d’évaluation.}
    \cpgeitem{Le niveau de maîtrise de ces compétences dépend de \textbf{l’autonomie et de l’initiative} requises dans les activités proposées aux étudiants sur les notions et capacités exigibles du programme.}
    \cpgeitem{La mise en œuvre des programmes doit aussi être l’occasion d’aborder avec les étudiants des questions liées à l’histoire de l’évolution des idées, des modèles et des théories en physique-chimie, à des questions liées à la recherche scientifique actuelle et à des enjeux citoyens comme la responsabilité individuelle et collective, la \textbf{sécurité} pour soi et pour autrui, l’\textbf{environnement} et le \textbf{développement durable} ou encore le \textbf{réchauffement climatique}.}},
}

\cpgerubrique{MP/P/3/APP}{
  parent = {MP/P/3},
  numero = {APP},
  niveau = {3},
  titre = {S’approprier},
  chapeau = {},
}

\cpgeligne{MP/P/3/APP/1}{
  parent = {MP/P/3/APP},
  numero = {APP.1},
  contenu = {
    \cpgeitem{S’approprier}},
  capacites = {
    \cpgeitem{Rechercher, extraire et organiser l’information en lien avec la situation étudiée.}},
  experimentales = {},
  numeriques = {},
}

\cpgeligne{MP/P/3/APP/2}{
  parent = {MP/P/3/APP},
  numero = {APP.2},
  contenu = {
    \cpgeitem{S’approprier}},
  capacites = {
    \cpgeitem{Identifier la complémentarité d’informations présentées sous des formes différentes (texte, graphe, tableau, etc.).}},
  experimentales = {},
  numeriques = {},
}

\cpgeligne{MP/P/3/APP/3}{
  parent = {MP/P/3/APP},
  numero = {APP.3},
  contenu = {
    \cpgeitem{S’approprier}},
  capacites = {
    \cpgeitem{Énoncer ou dégager une problématique scientifique.}},
  experimentales = {},
  numeriques = {},
}

\cpgeligne{MP/P/3/APP/4}{
  parent = {MP/P/3/APP},
  numero = {APP.4},
  contenu = {
    \cpgeitem{S’approprier}},
  capacites = {
    \cpgeitem{Représenter la situation par un schéma modèle.}},
  experimentales = {},
  numeriques = {},
}

\cpgeligne{MP/P/3/APP/5}{
  parent = {MP/P/3/APP},
  numero = {APP.5},
  contenu = {
    \cpgeitem{S’approprier}},
  capacites = {
    \cpgeitem{Identifier les grandeurs pertinentes, leur attribuer un symbole.}},
  experimentales = {},
  numeriques = {},
}

\cpgeligne{MP/P/3/APP/6}{
  parent = {MP/P/3/APP},
  numero = {APP.6},
  contenu = {
    \cpgeitem{S’approprier}},
  capacites = {
    \cpgeitem{Relier le problème à une situation modèle connue.}},
  experimentales = {},
  numeriques = {},
}

\cpgeligne{MP/P/3/APP/7}{
  parent = {MP/P/3/APP},
  numero = {APP.7},
  contenu = {
    \cpgeitem{S’approprier}},
  capacites = {
    \cpgeitem{Acquérir de nouvelles connaissances en autonomie.}},
  experimentales = {},
  numeriques = {},
}

\cpgerubrique{MP/P/3/ANA}{
  parent = {MP/P/3},
  numero = {ANA},
  niveau = {3},
  titre = {Analyser / Raisonner},
  chapeau = {},
}

\cpgeligne{MP/P/3/ANA/1}{
  parent = {MP/P/3/ANA},
  numero = {ANA.1},
  contenu = {
    \cpgeitem{Analyser / Raisonner}},
  capacites = {
    \cpgeitem{Formuler des hypothèses.}},
  experimentales = {},
  numeriques = {},
}

\cpgeligne{MP/P/3/ANA/2}{
  parent = {MP/P/3/ANA},
  numero = {ANA.2},
  contenu = {
    \cpgeitem{Analyser / Raisonner}},
  capacites = {
    \cpgeitem{Décomposer un problème en plusieurs problèmes plus simples.}},
  experimentales = {},
  numeriques = {},
}

\cpgeligne{MP/P/3/ANA/3}{
  parent = {MP/P/3/ANA},
  numero = {ANA.3},
  contenu = {
    \cpgeitem{Analyser / Raisonner}},
  capacites = {
    \cpgeitem{Proposer une stratégie pour répondre à une problématique.}},
  experimentales = {},
  numeriques = {},
}

\cpgeligne{MP/P/3/ANA/4}{
  parent = {MP/P/3/ANA},
  numero = {ANA.4},
  contenu = {
    \cpgeitem{Analyser / Raisonner}},
  capacites = {
    \cpgeitem{Choisir, concevoir, justifier un protocole, un dispositif expérimental, un modèle ou des lois physiques.}},
  experimentales = {},
  numeriques = {},
}

\cpgeligne{MP/P/3/ANA/5}{
  parent = {MP/P/3/ANA},
  numero = {ANA.5},
  contenu = {
    \cpgeitem{Analyser / Raisonner}},
  capacites = {
    \cpgeitem{Évaluer des ordres de grandeur.}},
  experimentales = {},
  numeriques = {},
}

\cpgeligne{MP/P/3/ANA/6}{
  parent = {MP/P/3/ANA},
  numero = {ANA.6},
  contenu = {
    \cpgeitem{Analyser / Raisonner}},
  capacites = {
    \cpgeitem{Identifier les idées essentielles d’un document et leurs articulations.}},
  experimentales = {},
  numeriques = {},
}

\cpgeligne{MP/P/3/ANA/7}{
  parent = {MP/P/3/ANA},
  numero = {ANA.7},
  contenu = {
    \cpgeitem{Analyser / Raisonner}},
  capacites = {
    \cpgeitem{Relier qualitativement ou quantitativement différents éléments d’un ou de documents.}},
  experimentales = {},
  numeriques = {},
}

\cpgerubrique{MP/P/3/REA}{
  parent = {MP/P/3},
  numero = {REA},
  niveau = {3},
  titre = {Réaliser},
  chapeau = {},
}

\cpgeligne{MP/P/3/REA/1}{
  parent = {MP/P/3/REA},
  numero = {REA.1},
  contenu = {
    \cpgeitem{Réaliser}},
  capacites = {
    \cpgeitem{Mettre en œuvre les étapes d’une démarche, un protocole, un modèle.}},
  experimentales = {},
  numeriques = {},
}

\cpgeligne{MP/P/3/REA/2}{
  parent = {MP/P/3/REA},
  numero = {REA.2},
  contenu = {
    \cpgeitem{Réaliser}},
  capacites = {
    \cpgeitem{Extraire une information d’un texte, d’un graphe, d’un tableau, d’un schéma, d’une photo.}},
  experimentales = {},
  numeriques = {},
}

\cpgeligne{MP/P/3/REA/3}{
  parent = {MP/P/3/REA},
  numero = {REA.3},
  contenu = {
    \cpgeitem{Réaliser}},
  capacites = {
    \cpgeitem{Schématiser un dispositif, une expérience, une méthode de mesure.}},
  experimentales = {},
  numeriques = {},
}

\cpgeligne{MP/P/3/REA/4}{
  parent = {MP/P/3/REA},
  numero = {REA.4},
  contenu = {
    \cpgeitem{Réaliser}},
  capacites = {
    \cpgeitem{Utiliser le matériel et les produits de manière adaptée en respectant des règles de sécurité.}},
  experimentales = {},
  numeriques = {},
}

\cpgeligne{MP/P/3/REA/5}{
  parent = {MP/P/3/REA},
  numero = {REA.5},
  contenu = {
    \cpgeitem{Réaliser}},
  capacites = {
    \cpgeitem{Effectuer des représentations graphiques à partir de données.}},
  experimentales = {},
  numeriques = {},
}

\cpgeligne{MP/P/3/REA/6}{
  parent = {MP/P/3/REA},
  numero = {REA.6},
  contenu = {
    \cpgeitem{Réaliser}},
  capacites = {
    \cpgeitem{Mener des calculs analytiques ou à l’aide d’un langage de programmation, effectuer des applications numériques.}},
  experimentales = {},
  numeriques = {},
}

\cpgeligne{MP/P/3/REA/7}{
  parent = {MP/P/3/REA},
  numero = {REA.7},
  contenu = {
    \cpgeitem{Réaliser}},
  capacites = {
    \cpgeitem{Conduire une analyse dimensionnelle.}},
  experimentales = {},
  numeriques = {},
}

\cpgerubrique{MP/P/3/VAL}{
  parent = {MP/P/3},
  numero = {VAL},
  niveau = {3},
  titre = {Valider},
  chapeau = {},
}

\cpgeligne{MP/P/3/VAL/1}{
  parent = {MP/P/3/VAL},
  numero = {VAL.1},
  contenu = {
    \cpgeitem{Valider}},
  capacites = {
    \cpgeitem{Exploiter des observations, des mesures en estimant les incertitudes.}},
  experimentales = {},
  numeriques = {},
}

\cpgeligne{MP/P/3/VAL/2}{
  parent = {MP/P/3/VAL},
  numero = {VAL.2},
  contenu = {
    \cpgeitem{Valider}},
  capacites = {
    \cpgeitem{Confronter les résultats d’un modèle à des résultats expérimentaux, à des données figurant dans un document, à ses connaissances.}},
  experimentales = {},
  numeriques = {},
}

\cpgeligne{MP/P/3/VAL/3}{
  parent = {MP/P/3/VAL},
  numero = {VAL.3},
  contenu = {
    \cpgeitem{Valider}},
  capacites = {
    \cpgeitem{Confirmer ou infirmer une hypothèse, une information.}},
  experimentales = {},
  numeriques = {},
}

\cpgeligne{MP/P/3/VAL/4}{
  parent = {MP/P/3/VAL},
  numero = {VAL.4},
  contenu = {
    \cpgeitem{Valider}},
  capacites = {
    \cpgeitem{Analyser les résultats de manière critique.}},
  experimentales = {},
  numeriques = {},
}

\cpgeligne{MP/P/3/VAL/5}{
  parent = {MP/P/3/VAL},
  numero = {VAL.5},
  contenu = {
    \cpgeitem{Valider}},
  capacites = {
    \cpgeitem{Repérer les points faibles d’une argumentation (contradiction, partialité, incomplétude, etc.).}},
  experimentales = {},
  numeriques = {},
}

\cpgeligne{MP/P/3/VAL/6}{
  parent = {MP/P/3/VAL},
  numero = {VAL.6},
  contenu = {
    \cpgeitem{Valider}},
  capacites = {
    \cpgeitem{Proposer des améliorations de la démarche ou du modèle.}},
  experimentales = {},
  numeriques = {},
}

\cpgerubrique{MP/P/3/COM}{
  parent = {MP/P/3},
  numero = {COM},
  niveau = {3},
  titre = {Communiquer},
  chapeau = {},
}

\cpgeligne{MP/P/3/COM/1}{
  parent = {MP/P/3/COM},
  numero = {COM.1},
  contenu = {
    \cpgeitem{Communiquer}},
  capacites = {
    \cpgeitem{À l’écrit comme à l’oral : \begin{itemize} \item présenter les étapes de sa démarche de manière synthétique, organisée et cohérente. \item rédiger une synthèse, une analyse, une argumentation. \item utiliser un vocabulaire scientifique précis et choisir des modes de représentation adaptés (schémas, graphes, cartes mentales, etc.). \end{itemize}}},
  experimentales = {},
  numeriques = {},
}

\cpgeligne{MP/P/3/COM/2}{
  parent = {MP/P/3/COM},
  numero = {COM.2},
  contenu = {
    \cpgeitem{Communiquer}},
  capacites = {
    \cpgeitem{Écouter, confronter son point de vue.}},
  experimentales = {},
  numeriques = {},
}

\cpgerubrique{MP/P/4}{
  parent = {MP/P},
  numero = {P.4},
  niveau = {2},
  titre = {Repères pour l’enseignement},
  chapeau = {
    \cpgeitem{Dans le cadre de la liberté pédagogique, l’enseignant organise son enseignement en respectant trois grands principes directeurs : \begin{itemize} \item privilégier la mise en activité des étudiants en évitant tout dogmatisme : l’acquisition des connaissances, des capacités et des compétences est d’autant plus efficace que les étudiants sont acteurs de leur formation. Les supports pédagogiques utilisés doivent notamment favoriser la réflexion, le raisonnement, la participation et l’autonomie des étudiants. L’investigation expérimentale et la résolution de problèmes facilitent cette mise en activité ; \item recourir à la mise en contexte des contenus scientifiques : le questionnement scientifique peut être introduit à partir de phénomènes naturels, de procédés industriels ou d’objets technologiques. Le recours à des approches documentaires est un moyen pertinent pour diversifier les supports d’accès à l’information scientifique et technologique et ainsi former l’étudiant à mieux en appréhender la complexité et à apprendre par lui-même. Lorsque le thème traité s’y prête, l’enseignant peut le mettre en perspective avec l’histoire des sciences et des techniques, avec des questions d’actualité ou des débats d’idées ; \item contribuer à la nécessaire mise en cohérence des enseignements scientifiques ; la progression en physique-chimie doit être articulée avec celles mises en œuvre dans les autres disciplines scientifiques : mathématiques, informatique, sciences industrielles de l’ingénieur. \end{itemize}}
    \cpgeitem{Concernant l’évaluation, qui vise à mesurer le degré de maîtrise du socle ainsi défini et le niveau d’autonomie et d’initiative des étudiants, l’enseignant veille soigneusement à identifier les compétences et les capacités mobilisées dans les activités proposées afin d’en élargir le plus possible le spectre.}
    \cpgeitem{Enfin, le professeur veille aussi à développer chez les étudiants des compétences transversales et préprofessionnelles relatives aux capacités suivantes : \begin{itemize} \item identifier les différents champs professionnels et les parcours pour y accéder ; \item valoriser ses compétences scientifiques et techniques en lien avec son projet de poursuite d'études ou professionnel. \end{itemize}}},
}

\cpgerubrique{MP/E}{
  parent = {MP},
  numero = {E},
  niveau = {1},
  titre = {Formation expérimentale},
  chapeau = {
    \cpgeitem{Cette partie est spécifiquement dédiée à la mise en œuvre de la formation expérimentale des étudiants lors des séances de travaux pratiques.}
    \cpgeitem{Dans un premier temps, elle précise les connaissances et savoir-faire qui doivent être acquis dans le domaine de la mesure et de l’évaluation des incertitudes. Elle présente ensuite de façon détaillée l’ensemble des capacités expérimentales qui doivent être acquises en autonomie par les étudiants à l’issue de leur seconde année de CPGE. Enfin, elle aborde la question de la prévention du risque au laboratoire de physique-chimie.}
    \cpgeitem{Une liste de matériel, que les étudiants doivent savoir utiliser avec l’aide d’une notice succincte, figure dans l’annexe 1 du présent programme.}},
}

\cpgerubrique{MP/E/1}{
  parent = {MP/E},
  numero = {E.1},
  niveau = {2},
  titre = {Mesures et incertitudes},
  chapeau = {
    \cpgeitem{Les notions et capacités identifiées ci-dessous couvrent les deux années de formation en classe préparatoire aux grandes écoles ; leur pleine maîtrise est donc un objectif de fin de seconde année.}
    \cpgeitem{L’accent est mis sur la variabilité de la mesure d’une grandeur physique et sa caractérisation à l’aide de l’incertitude-type. La comparaison entre deux valeurs mesurées d’une même grandeur physique est conduite au moyen de l’écart normalisé, l’objectif principal étant de développer l’esprit critique des étudiants en s’appuyant sur un critère quantitatif. Le même esprit prévaut dans l’analyse des résultats d’une régression linéaire qui ne saurait s’appuyer sur l’exploitation non raisonnée du coefficient de corrélation ($R^2$).}
    \cpgeitem{Le recours à la simulation vise à illustrer, sur la base de mesures expérimentales, différents effets de la variabilité de la mesure d’une grandeur physique dans les cas des incertitudes-types composées et de la régression linéaire.}},
}

\cpgeligne{MP/E/1/1}{
  parent = {MP/E/1},
  numero = {E.1.1},
  contenu = {
    \cpgeitem{Variabilité de la mesure d’une grandeur physique. Incertitude.}
    \cpgeitem{Incertitude-type.}},
  capacites = {
    \cpgeitem{Identifier les incertitudes liées, par exemple, à l’opérateur, à l’environnement, aux instruments ou à la méthode de mesure.}
    \cpgeitem{Procéder à l’évaluation d’une incertitude-type par une approche statistique (évaluation de type A).}
    \cpgeitem{Procéder à l’évaluation d’une incertitude-type par une autre approche que statistique (évaluation de type B).}
    \cpgeitem{Associer un intervalle de confiance à l’écart-type dans l’hypothèse d’une distribution suivant la loi normale.}},
  experimentales = {},
  numeriques = {},
}

\cpgeligne{MP/E/1/2}{
  parent = {MP/E/1},
  numero = {E.1.2},
  contenu = {
    \cpgeitem{Incertitudes-types composées.}},
  capacites = {
    \cpgeitem{Évaluer l’incertitude-type d’une grandeur s’exprimant en fonction d’autres grandeurs, dont les incertitudes-types sont connues, à l’aide d’une somme, d’une différence, d’un produit ou d’un quotient.}
    \cpgeitem{Comparer entre elles les différentes contributions lors de l’évaluation d’une incertitude-type composée.}},
  experimentales = {},
  numeriques = {
    \cpgeitem{simuler, à l’aide d’un langage de programmation ou d’un tableur, un processus aléatoire permettant de caractériser la variabilité de la valeur d’une grandeur composée.}},
}

\cpgeligne{MP/E/1/3}{
  parent = {MP/E/1},
  numero = {E.1.3},
  contenu = {
    \cpgeitem{Écriture du résultat d’une mesure.}},
  capacites = {
    \cpgeitem{Écrire, avec un nombre adapté de chiffres significatifs, le résultat d’une mesure.}},
  experimentales = {},
  numeriques = {},
}

\cpgeligne{MP/E/1/4}{
  parent = {MP/E/1},
  numero = {E.1.4},
  contenu = {
    \cpgeitem{Comparaison de deux valeurs ; écart normalisé.}},
  capacites = {
    \cpgeitem{Comparer deux valeurs dont les incertitudes-types sont connues à l’aide de leur écart normalisé.}
    \cpgeitem{Analyser les causes d’une éventuelle incompatibilité entre le résultat d’une mesure et le résultat attendu par une modélisation.}},
  experimentales = {},
  numeriques = {},
}

\cpgeligne{MP/E/1/5}{
  parent = {MP/E/1},
  numero = {E.1.5},
  contenu = {
    \cpgeitem{Régression linéaire.}},
  capacites = {
    \cpgeitem{Utiliser un logiciel de régression linéaire afin d’obtenir les valeurs des paramètres du modèle.}
    \cpgeitem{Analyser les résultats obtenus à l’aide d’une procédure de validation : analyse graphique intégrant les barres d’incertitude ou analyse des écarts normalisés.}},
  experimentales = {},
  numeriques = {
    \cpgeitem{simuler, à l’aide d’un langage de programmation ou d’un tableur, un processus aléatoire de variation des valeurs expérimentales de l’une des grandeurs – simulation Monte-Carlo – pour évaluer l’incertitude sur les paramètres du modèle.}},
}

\cpgerubrique{MP/E/2}{
  parent = {MP/E},
  numero = {E.2},
  niveau = {2},
  titre = {Mesures et capacités expérimentales},
  chapeau = {
    \cpgeitem{Cette partie présente l’ensemble des capacités expérimentales nouvelles que les étudiants doivent acquérir au cours de l’année de MP durant les séances de travaux pratiques. Elle vient   prolonger la partie correspondante du programme de MPSI dont les capacités doivent être complètement acquises à l’issue des deux années de préparation, et restent donc au programme de la classe de MP.}
    \cpgeitem{Les capacités rassemblées ici ne constituent en aucun cas une liste de travaux pratiques qui s’articuleraient autour d’une découverte du matériel, mais doivent au contraire faire l’objet d’un apprentissage progressif contextualisé où chaque élément apparaît naturellement à l’occasion d’un problème concret. À ce titre, elle vient compléter la liste des thèmes d’étude – en gras dans la colonne « capacités exigibles » de la partie « Contenus thématiques » – à partir desquels la problématique d’une séance peut être définie.}
    \cpgeitem{Les activités expérimentales sur le thème de la chimie sont aussi l’occasion de consolider les savoir-faire de la classe de MPSI en particulier dans le domaine des solutions aqueuses.}},
}

\cpgeligne{MP/E/2/1}{
  parent = {MP/E/2},
  numero = {E.2.1},
  contenu = {
    \cpgeitem{Analyse spectrale.}},
  capacites = {
    \cpgeitem{Mettre en évidence le phénomène de repliement du spectre provoqué par l’échantillonnage avec un oscilloscope numérique ou une carte d’acquisition.\cpgerenvoi{MP/2/2/1}}
    \cpgeitem{Choisir les paramètres d’une acquisition numérique destinée à une analyse spectrale afin de respecter la condition de Nyquist-Shannon, tout en optimisant la résolution spectrale.\cpgerenvoi{MP/2/2/1}}},
  experimentales = {},
  numeriques = {},
}

\cpgeligne{MP/E/2/2}{
  parent = {MP/E/2},
  numero = {E.2.2},
  contenu = {
    \cpgeitem{Filtrage analogique d’un signal périodique.}},
  capacites = {
    \cpgeitem{Mettre en évidence l’action d’un filtre linéaire sur un signal périodique dans les domaines fréquentiel et temporel.\cpgerenvoi{MP/2/1/2}}},
  experimentales = {},
  numeriques = {},
}

\cpgeligne{MP/E/2/3}{
  parent = {MP/E/2},
  numero = {E.2.3},
  contenu = {
    \cpgeitem{Montages utilisant un amplificateur linéaire intégré (ALI).}},
  capacites = {
    \cpgeitem{Mettre en œuvre divers montages utilisant un ALI, les schémas des montages étant fournis.}},
  experimentales = {},
  numeriques = {},
}

\cpgeligne{MP/E/2/4}{
  parent = {MP/E/2},
  numero = {E.2.4},
  contenu = {
    \cpgeitem{Électronique numérique.}},
  capacites = {
    \cpgeitem{Utiliser un convertisseur analogique-numérique et un convertisseur numérique-analogique.\cpgerenvoi{MP/2/2/2}}},
  experimentales = {},
  numeriques = {},
}

\cpgeligne{MP/E/2/5}{
  parent = {MP/E/2},
  numero = {E.2.5},
  contenu = {
    \cpgeitem{Ondes électromagnétiques.}},
  capacites = {
    \cpgeitem{Mettre en œuvre un détecteur d’ondes électromagnétiques dans le domaine des ondes centimétriques.\cpgerenvoi{MP/4/5/8}}},
  experimentales = {},
  numeriques = {},
}

\cpgeligne{MP/E/2/6}{
  parent = {MP/E/2},
  numero = {E.2.6},
  contenu = {
    \cpgeitem{Analyse d’une lumière.}},
  capacites = {
    \cpgeitem{Identifier, à l’aide d’un polariseur, une onde polarisée rectilignement et déterminer sa direction de polarisation.\cpgerenvoi{MP/4/5/3}}
    \cpgeitem{Mesurer une longueur d'onde à l’aide d’un goniomètre équipé d'un réseau.\cpgerenvoi{MP/3/2/3}}},
  experimentales = {},
  numeriques = {},
}

\cpgeligne{MP/E/2/7}{
  parent = {MP/E/2},
  numero = {E.2.7},
  contenu = {
    \cpgeitem{Analyse d’une figure d'interférence.}},
  capacites = {
    \cpgeitem{Mettre en œuvre un photodétecteur en sortie d’un interféromètre.}},
  experimentales = {},
  numeriques = {},
}

\cpgeligne{MP/E/2/8}{
  parent = {MP/E/2},
  numero = {E.2.8},
  contenu = {
    \cpgeitem{Étude de la cohérence temporelle d’une source.}},
  capacites = {
    \cpgeitem{Régler un interféromètre de Michelson compensé pour une observation en lame d’air avec une source étendue à l’aide d’un protocole fourni.\cpgerenvoi{MP/3/4/2}}
    \cpgeitem{Obtenir une estimation de la longueur de cohérence d’une source et de l’écart spectral d’un doublet à l’aide d’un interféromètre de Michelson réglé en lame d’air.\cpgerenvoi{MP/3/4/2}}},
  experimentales = {},
  numeriques = {},
}

\cpgeligne{MP/E/2/9}{
  parent = {MP/E/2},
  numero = {E.2.9},
  contenu = {
    \cpgeitem{Conduction thermique.}},
  capacites = {
    \cpgeitem{Mettre en œuvre un dispositif de mesure de conductivité thermique.\cpgerenvoi{MP/5/2/4}}},
  experimentales = {},
  numeriques = {},
}

\cpgeligne{MP/E/2/10}{
  parent = {MP/E/2},
  numero = {E.2.10},
  contenu = {
    \cpgeitem{Bilans d’énergie.}},
  capacites = {
    \cpgeitem{Mettre en œuvre une technique de calorimétrie.}
    \cpgeitem{Déterminer la valeur en eau d’un calorimètre.}
    \cpgeitem{Estimer les fuites thermiques lors d’expériences réalisées avec un calorimètre.}},
  experimentales = {},
  numeriques = {},
}

\cpgeligne{MP/E/2/11}{
  parent = {MP/E/2},
  numero = {E.2.11},
  contenu = {
    \cpgeitem{Mesures de grandeurs électriques. : conductance-conductivité, tension électrique, intensité du courant.}},
  capacites = {
    \cpgeitem{Mettre en œuvre des mesures de grandeurs physiques pour déterminer la composition d’un système chimique.}},
  experimentales = {},
  numeriques = {},
}

\cpgeligne{MP/E/2/12}{
  parent = {MP/E/2},
  numero = {E.2.12},
  contenu = {
    \cpgeitem{Électrochimie.}},
  capacites = {
    \cpgeitem{Mettre en œuvre un dispositif à trois électrodes pour tracer des courbes courant-potentiel.\cpgerenvoi{MP/9/2/1}}
    \cpgeitem{Mettre en œuvre des piles et des électrolyseurs. \cpgerenvoi{MP/9/3/4}}},
  experimentales = {},
  numeriques = {},
}

\cpgerubrique{MP/E/3}{
  parent = {MP/E},
  numero = {E.3},
  niveau = {2},
  titre = {Prévention du risque au laboratoire de physique-chimie},
  chapeau = {
    \cpgeitem{Les étudiants doivent prendre conscience du risque lié à la manipulation et au rejet des produits chimiques. L’apprentissage et le respect des règles de sécurité chimique, électrique, optique et celles liées à la pression et à la température leur permettent de prévenir et de minimiser ce risque.}
    \cpgeitem{Futurs ingénieurs, chercheurs, enseignants, ils doivent être sensibilisés au respect de la législation et à l’impact de leur activité sur l’environnement.}},
}

\cpgeligne{MP/E/3/1}{
  parent = {MP/E/3},
  numero = {E.3.1},
  contenu = {
    \cpgeitem{\cpgenotion{1. Prévention des risques au laboratoire}}},
  capacites = {
    \cpgeitem{Adopter une attitude responsable et adaptée au travail en laboratoire.}
    \cpgeitem{Développer une attitude autonome dans la prévention des risques.}},
  experimentales = {},
  numeriques = {},
}

\cpgeligne{MP/E/3/2}{
  parent = {MP/E/3},
  numero = {E.3.2},
  contenu = {
    \cpgeitem{\cpgenotion{Risque chimique}}
    \cpgeitem{Règles de sécurité au laboratoire.}
    \cpgeitem{Classes et catégories de danger.}
    \cpgeitem{Pictogrammes de sécurité pour les produits chimiques.}
    \cpgeitem{Mentions de danger (H) et conseils de prudence (P).}
    \cpgeitem{Fiches de sécurité.}},
  capacites = {
    \cpgeitem{Relever les indications sur le risque associé au prélèvement, au mélange et au stockage des produits chimiques et adopter une attitude responsable lors de leur utilisation.}},
  experimentales = {},
  numeriques = {},
}

\cpgeligne{MP/E/3/3}{
  parent = {MP/E/3},
  numero = {E.3.3},
  contenu = {
    \cpgeitem{Risque électrique}},
  capacites = {
    \cpgeitem{Adopter une attitude responsable lors de l’utilisation d’appareils électriques.}},
  experimentales = {},
  numeriques = {},
}

\cpgeligne{MP/E/3/4}{
  parent = {MP/E/3},
  numero = {E.3.4},
  contenu = {
    \cpgeitem{Risque optique}},
  capacites = {
    \cpgeitem{Utiliser les sources laser et les diodes électroluminescentes de manière adaptée.}},
  experimentales = {},
  numeriques = {},
}

\cpgeligne{MP/E/3/5}{
  parent = {MP/E/3},
  numero = {E.3.5},
  contenu = {
    \cpgeitem{Risques liés à la pression et à la température}},
  capacites = {
    \cpgeitem{Adopter une attitude responsable lors de manipulations de corps chauds ou de dispositifs engageant des hautes ou des}
    \cpgeitem{basses pressions.}},
  experimentales = {},
  numeriques = {},
}

\cpgeligne{MP/E/3/6}{
  parent = {MP/E/3},
  numero = {E.3.6},
  contenu = {
    \cpgeitem{\cpgenotion{2. Prévention de l’impact environnemental}}
    \cpgeitem{Traitement et rejet des espèces chimiques.}},
  capacites = {
    \cpgeitem{Adapter le mode d’élimination d’une espèce chimique ou d’un mélange en fonction des informations recueillies sur la toxicité ou les risques.}
    \cpgeitem{Sélectionner, parmi plusieurs modes opératoires, celui qui minimise les impacts environnementaux.}},
  experimentales = {},
  numeriques = {},
}

\cpgerubrique{MP/1}{
  parent = {MP},
  numero = {1},
  niveau = {1},
  titre = {Mécanique},
  chapeau = {
    \cpgeitem{Le programme de mécanique de MP vise à compléter les acquis de mécanique du cours de MPSI. Il est structuré en deux sous-parties : la première est consacrée aux changements de référentiels, la seconde aux conséquences mécaniques des actions de frottements entre solides.}},
}

\cpgerubrique{MP/1/1}{
  parent = {MP/1},
  numero = {1.1},
  niveau = {2},
  titre = {Référentiels non galiléens},
  chapeau = {
    \cpgeitem{La partie intitulée \cpgecettepartie{} est organisée autour de deux situations : la translation et la rotation uniforme autour d’un axe fixe. L’accent est mis sur la compréhension qualitative des effets observés, l’évaluation des ordres de grandeurs et les conséquences expérimentales.}},
}

\cpgeligne{MP/1/1/1}{
  parent = {MP/1/1},
  numero = {1.1.1},
  contenu = {
    \cpgeitem{Mouvement d’un référentiel par rapport à un autre dans les cas du mouvement de translation et du mouvement de rotation uniforme autour d’un axe fixe.}},
  capacites = {
    \cpgeitem{Reconnaître et caractériser un mouvement de translation et un mouvement de rotation uniforme autour d’un axe fixe d’un référentiel par rapport à un autre.}},
  experimentales = {},
  numeriques = {},
}

\cpgeligne{MP/1/1/2}{
  parent = {MP/1/1},
  numero = {1.1.2},
  contenu = {
    \cpgeitem{Vecteur rotation d'un référentiel par rapport à un autre.}},
  capacites = {
    \cpgeitem{Exprimer le vecteur rotation d'un référentiel par rapport à un autre.}},
  experimentales = {},
  numeriques = {},
}

\cpgeligne{MP/1/1/3}{
  parent = {MP/1/1},
  numero = {1.1.3},
  contenu = {
    \cpgeitem{Composition des vitesses et des accélérations dans le cas d’une translation, et dans le cas d’une rotation uniforme autour d’un axe fixe : vitesse d’entraînement, accélérations d’entraînement et de Coriolis.}},
  capacites = {
    \cpgeitem{Relier les dérivées d'un vecteur dans des référentiels différents par la relation de la dérivation composée.}
    \cpgeitem{Citer et utiliser les expressions de la vitesse d’entraînement et des accélérations d’entraînement et de Coriolis.}},
  experimentales = {},
  numeriques = {},
}

\cpgeligne{MP/1/1/4}{
  parent = {MP/1/1},
  numero = {1.1.4},
  contenu = {
    \cpgeitem{Dynamique du point en référentiel non galiléen dans le cas où le référentiel entraîné est en translation, ou en rotation uniforme autour d’un axe fixe par rapport à un référentiel galiléen. Forces d’inertie.}},
  capacites = {
    \cpgeitem{Exprimer les forces d'inerties, dans les seuls cas où le référentiel entraîné est en translation, ou en rotation uniforme autour d’un axe fixe par rapport à un référentiel galiléen.}
    \cpgeitem{Décrire et interpréter les effets des forces d’inertie dans des cas concrets : sens de la force d’inertie d’entraînement dans un mouvement de translation ; caractère centrifuge de la force d'inertie d’entraînement dans le cas où le référentiel est en rotation uniforme autour d’un axe fixe par rapport à un référentiel galiléen.}
    \cpgeitem{Utiliser les lois de la dynamique en référentiel non galiléen dans les seuls cas où le référentiel entraîné est en translation, ou en rotation uniforme autour d’un axe fixe par rapport à un référentiel galiléen.}},
  experimentales = {},
  numeriques = {},
}

\cpgeligne{MP/1/1/5}{
  parent = {MP/1/1},
  numero = {1.1.5},
  contenu = {
    \cpgeitem{Caractère galiléen approché d’un référentiel. Exemple du référentiel de Copernic, du référentiel géocentrique et du référentiel terrestre.}},
  capacites = {
    \cpgeitem{Citer quelques manifestations du caractère non galiléen du référentiel terrestre.}
    \cpgeitem{Estimer, en ordre de grandeur, la contribution de la force d’inertie de Coriolis dans un problème de dynamique terrestre.}},
  experimentales = {},
  numeriques = {},
}

\cpgerubrique{MP/1/2}{
  parent = {MP/1},
  numero = {1.2},
  niveau = {2},
  titre = {Lois du frottement solide},
  chapeau = {
    \cpgeitem{La partie \cpgecettepartie{}, est limitée au seul cas de la translation ; elle permet de mettre en œuvre un mode de raisonnement spécifique et particulièrement formateur, sans pour autant omettre les conséquences expérimentales.}},
}

\cpgeligne{MP/1/2/1}{
  parent = {MP/1/2},
  numero = {1.2.1},
  contenu = {
    \cpgeitem{Contact entre deux solides. Aspects microscopiques.}
    \cpgeitem{Lois de Coulomb du frottement de glissement dans le seul cas d’un solide en translation.}
    \cpgeitem{Aspect énergétique.}},
  capacites = {
    \cpgeitem{Utiliser les lois de Coulomb dans les trois situations : équilibre, mise en mouvement, freinage.}
    \cpgeitem{Formuler une hypothèse (quant au glissement ou non) et la valider.}
    \cpgeitem{Effectuer un bilan énergétique.}},
  experimentales = {
    \cpgeitem{Effectuer une mesure d’un coefficient de frottement.}},
  numeriques = {
    \cpgeitem{à l’aide d’un langage de programmation, simuler une situation mécanique dans laquelle intervient au moins un changement de mode de glissement}},
}

\cpgerubrique{MP/2}{
  parent = {MP},
  numero = {2},
  niveau = {1},
  titre = {Éléments de traitement du signal},
  chapeau = {
    \cpgeitem{Ce thème du programme, décomposé en deux parties, complète l’étude des circuits électriques linéaires menée dans la partie « Ondes et signaux » du programme de MPSI. La composante expérimentale est forte et les capacités exigibles ont vocation à être principalement développées au cours de séances de travaux pratiques.}},
}

\cpgerubrique{MP/2/1}{
  parent = {MP/2},
  numero = {2.1},
  niveau = {2},
  titre = {Signaux périodiques},
  chapeau = {
    \cpgeitem{Dans la première partie intitulée \cpgecettepartie{}, l’accent est mis sur l’action d’un filtre linéaire sur un signal périodique, l’objectif étant de comprendre le rôle central de la linéarité des systèmes pour interpréter ou prévoir la forme du signal résultant d’un filtrage.}},
}

\cpgeligne{MP/2/1/1}{
  parent = {MP/2/1},
  numero = {2.1.1},
  contenu = {
    \cpgeitem{Signaux périodiques.}},
  capacites = {
    \cpgeitem{Commenter le spectre d'un signal périodique : relier la décomposition spectrale et l’allure du signal dans le domaine temporel.}},
  experimentales = {},
  numeriques = {},
}

\cpgeligne{MP/2/1/2}{
  parent = {MP/2/1},
  numero = {2.1.2},
  contenu = {
    \cpgeitem{Action d’un filtre linéaire du premier ou du second ordre sur un signal périodique.}},
  capacites = {
    \cpgeitem{Prévoir l’effet d’un filtrage linéaire sur la composition spectrale d’un signal périodique. Expliciter les conditions pour obtenir un comportement intégrateur ou dérivateur.}},
  experimentales = {
    \cpgeitem{Mettre en œuvre un dispositif expérimental illustrant l’action d’un filtre sur un signal périodique.\cpgerenvoi{MP/E/2/2}}},
  numeriques = {},
}

\cpgerubrique{MP/2/2}{
  parent = {MP/2},
  numero = {2.2},
  niveau = {2},
  titre = {Électronique numérique},
  chapeau = {
    \cpgeitem{La seconde partie intitulée \cpgecettepartie{} est à vocation uniquement expérimentale ; elle constitue une initiation au traitement numérique des signaux à travers les points suivants : l’échantillonnage et le repliement de spectre, la conversion analogique/numérique et le filtrage numérique. Le phénomène de repliement de spectre est présenté qualitativement au moyen d’illustrations démonstratives, l’objectif étant de mettre en place la condition de Nyquist-Shannon afin de réaliser convenablement une acquisition numérique.}},
}

\cpgeligne{MP/2/2/1}{
  parent = {MP/2/2},
  numero = {2.2.1},
  contenu = {
    \cpgeitem{Échantillonnage, fréquence d’échantillonnage. Conséquences expérimentales du théorème de Nyquist-Shannon.}},
  capacites = {},
  experimentales = {
    \cpgeitem{Réaliser l'échantillonnage d'un signal.}
    \cpgeitem{Choisir la fréquence d’échantillonnage afin de respecter la condition de Nyquist-Shannon.\cpgerenvoi{MP/E/2/1}}
    \cpgeitem{Commenter la structure du spectre du signal obtenu après échantillonnage.}
    \cpgeitem{Mettre en évidence le phénomène de repliement de spectre au moyen d'un}
    \cpgeitem{oscilloscope numérique ou d'un logiciel de calcul numérique.\cpgerenvoi{MP/E/2/1}}},
  numeriques = {},
}

\cpgeligne{MP/2/2/2}{
  parent = {MP/2/2},
  numero = {2.2.2},
  contenu = {
    \cpgeitem{Filtrage numérique.}},
  capacites = {},
  experimentales = {
    \cpgeitem{Mettre en œuvre un convertisseur analogique/numérique et un traitement numérique afin de réaliser un filtre passe-bas ; utiliser un convertisseur}
    \cpgeitem{numérique/analogique pour restituer un signal analogique.\cpgerenvoi{E.2.4}}},
  numeriques = {
    \cpgeitem{à l’aide d’un langage de programmation, simuler un filtrage numérique et visualiser son action sur un signal périodique.}},
}

\cpgerubrique{MP/3}{
  parent = {MP},
  numero = {3},
  niveau = {1},
  titre = {Optique},
  chapeau = {
    \cpgeitem{Le programme d’optique de la filière MP s’inscrit dans le prolongement du thème « Ondes et signaux » du programme de MPSI. Il s’agit pour les étudiants d'approfondir l'étude des phénomènes d'interférences lumineuses, dans le cadre du modèle ondulatoire de la lumière. L’approche reste centrée sur l’expérience, mais la modélisation doit permettre d’analyser de façon raisonnée les conditions optimales d’observation d’interférences lumineuses, et leur exploitation quantitative. L'enseignant ne manquera pas de rappeler que ces phénomènes, étudiés ici dans le cadre de l'optique, sont généralisables à tout comportement ondulatoire.}},
}

\cpgerubrique{MP/3/1}{
  parent = {MP/3},
  numero = {3.1},
  niveau = {2},
  titre = {Modèle scalaire des ondes lumineuses},
  chapeau = {
    \cpgeitem{La partie \cpgecettepartie{} introduit les outils nécessaires pour décrire les phénomènes d’interférences. Le programme utilise le mot « intensité » pour décrire la grandeur détectée mais on peut utiliser indifféremment les mots « intensité » ou « éclairement » sans chercher à les distinguer à ce niveau. L’intensité lumineuse est introduite comme une puissance par unité de surface. Le théorème de Malus (orthogonalité des rayons de lumière et des surfaces d’ondes) est admis.}},
}

\cpgeligne{MP/3/1/1}{
  parent = {MP/3/1},
  numero = {3.1.1},
  contenu = {
    \cpgeitem{Modèle de propagation dans l’approximation de l’optique géométrique.}},
  capacites = {
    \cpgeitem{Utiliser une grandeur scalaire pour décrire un signal lumineux.}},
  experimentales = {},
  numeriques = {},
}

\cpgeligne{MP/3/1/2}{
  parent = {MP/3/1},
  numero = {3.1.2},
  contenu = {
    \cpgeitem{Chemin optique. Déphasage dû à la propagation. Surfaces d’ondes. Théorème de Malus (admis).}
    \cpgeitem{Onde plane, onde sphérique ; effet d’une lentille mince dans l’approximation de Gauss.}},
  capacites = {
    \cpgeitem{Exprimer le retard de phase en un point (par rapport à un autre) en fonction de la durée de propagation ou du chemin optique.}
    \cpgeitem{Associer une description de la formation des images en termes de rayon lumineux et en termes de surfaces d’onde.}
    \cpgeitem{Utiliser la propriété énonçant que le chemin optique séparant deux points conjugués est indépendant du rayon lumineux choisi.}},
  experimentales = {},
  numeriques = {},
}

\cpgeligne{MP/3/1/3}{
  parent = {MP/3/1},
  numero = {3.1.3},
  contenu = {
    \cpgeitem{Modèle d’émission. Relation (admise) entre le temps de cohérence et la largeur spectrale.}},
  capacites = {
    \cpgeitem{Citer l'ordre de grandeur du temps de cohérence $\Delta t$ de quelques radiations visibles.}
    \cpgeitem{Utiliser la relation $\Delta f.\Delta t\sim1$ pour relier le temps de cohérence à la largeur spectrale $\Delta\lambda$ de la radiation.}},
  experimentales = {},
  numeriques = {},
}

\cpgeligne{MP/3/1/4}{
  parent = {MP/3/1},
  numero = {3.1.4},
  contenu = {
    \cpgeitem{Récepteurs. Intensité de la lumière.}},
  capacites = {
    \cpgeitem{Relier l’intensité à la moyenne temporelle du carré de la grandeur scalaire de l’optique.}
    \cpgeitem{Citer l’ordre de grandeur du temps de réponse de quelques récepteurs de lumière.}},
  experimentales = {
    \cpgeitem{Mettre en œuvre des expériences utilisant un capteur photographique numérique.}},
  numeriques = {},
}

\cpgerubrique{MP/3/2}{
  parent = {MP/3},
  numero = {3.2},
  niveau = {2},
  titre = {Superposition d’ondes lumineuses},
  chapeau = {
    \cpgeitem{Dans la partie \cpgecettepartie{}, la formule de Fresnel, admise en classe de première année, est démontrée. L’étude de la superposition de N ondes cohérentes ne doit pas donner lieu à des développements calculatoires.}},
}

\cpgeligne{MP/3/2/1}{
  parent = {MP/3/2},
  numero = {3.2.1},
  contenu = {
    \cpgeitem{Superposition de deux ondes incohérentes entre elles.}},
  capacites = {
    \cpgeitem{Justifier et utiliser l’additivité des intensités.}},
  experimentales = {},
  numeriques = {},
}

\cpgeligne{MP/3/2/2}{
  parent = {MP/3/2},
  numero = {3.2.2},
  contenu = {
    \cpgeitem{Superposition de deux ondes monochromatiques cohérentes entre elles : formule de Fresnel. Facteur de contraste.}},
  capacites = {
    \cpgeitem{Citer les principales conditions pour que le phénomène d'interférences apparaisse (ondes quasi synchrones, déphasage constant dans le temps ou très lentement variable).}
    \cpgeitem{Établir et utiliser la formule de Fresnel. Associer un bon contraste à des ondes d’intensités voisines.}},
  experimentales = {},
  numeriques = {},
}

\cpgeligne{MP/3/2/3}{
  parent = {MP/3/2},
  numero = {3.2.3},
  contenu = {
    \cpgeitem{Superposition de N ondes monochromatiques cohérentes entre elles, de même amplitude et dont les phases sont en progression arithmétique.}},
  capacites = {
    \cpgeitem{Établir la relation fondamentale des réseaux liant la condition d'interférences constructives à l’expression de la différence de marche entre deux ondes issues de motifs consécutifs.}
    \cpgeitem{Établir, par le calcul, la demi-largeur $2\pi/N$ des pics principaux de la courbe d’intensité en fonction du déphasage.}},
  experimentales = {
    \cpgeitem{Mettre en œuvre un dispositif expérimental utilisant un phénomène d’interférences à N ondes.\cpgerenvoi{MP/E/2/6}}},
  numeriques = {},
}

\cpgerubrique{MP/3/3}{
  parent = {MP/3},
  numero = {3.3},
  niveau = {2},
  titre = {Exemple de dispositif interférentiel par division du front d’onde : trous d’Young},
  chapeau = {
    \cpgeitem{Dans la partie \cpgecettepartie{}, les trous d’Young permettent de confronter théorie et expérience. Les fentes d’Young peuvent être abordées mais de manière exclusivement expérimentale. Aucune connaissance sur un autre diviseur du front d’onde n’est exigible.}},
}

\cpgeligne{MP/3/3/1}{
  parent = {MP/3/3},
  numero = {3.3.1},
  contenu = {
    \cpgeitem{Trous d’Young ponctuels dans un milieu non dispersif : source ponctuelle à distance finie et observation à grande distance. Champ}
    \cpgeitem{d'interférences. Ordre d’interférences.}},
  capacites = {
    \cpgeitem{Définir, exprimer et utiliser l’interfrange et l’ordre d’interférences.}
    \cpgeitem{Justifier que les franges ne sont pas localisées.}},
  experimentales = {},
  numeriques = {},
}

\cpgeligne{MP/3/3/2}{
  parent = {MP/3/3},
  numero = {3.3.2},
  contenu = {
    \cpgeitem{Variations de l’ordre d’interférences avec la position du point d’observation ; franges d’interférences.}},
  capacites = {
    \cpgeitem{Interpréter la forme des franges observées.}},
  experimentales = {},
  numeriques = {},
}

\cpgeligne{MP/3/3/3}{
  parent = {MP/3/3},
  numero = {3.3.3},
  contenu = {
    \cpgeitem{Variations de l’ordre d’interférences avec la position d’un point source.}
    \cpgeitem{Perte de contraste par élargissement angulaire de la source.}},
  capacites = {
    \cpgeitem{Utiliser un critère de brouillage des franges portant sur l’ordre d’interférence.}},
  experimentales = {},
  numeriques = {},
}

\cpgeligne{MP/3/3/4}{
  parent = {MP/3/3},
  numero = {3.3.4},
  contenu = {
    \cpgeitem{Variations de l’ordre d’interférence avec la longueur d’onde.}
    \cpgeitem{Perte de contraste par élargissement spectral de la source.}},
  capacites = {
    \cpgeitem{Utiliser un critère de brouillage des franges portant sur l’ordre d’interférence.}},
  experimentales = {},
  numeriques = {},
}

\cpgerubrique{MP/3/4}{
  parent = {MP/3},
  numero = {3.4},
  niveau = {2},
  titre = {Exemple de dispositif interférentiel par division d’amplitude : interféromètre de Michelson éclairé par une source spatialement étendue},
  chapeau = {
    \cpgeitem{Dans la partie \cpgecettepartie{}, l’étude de l’interféromètre de Michelson en lame d’air permet de confronter théorie et expérience. L’étude de l’interféromètre de Michelson en coin d’air est abordée de manière exclusivement expérimentale. Pour la modélisation d’un interféromètre de Michelson, on suppose la séparatrice infiniment mince.}},
}

\cpgeligne{MP/3/4/1}{
  parent = {MP/3/4},
  numero = {3.4.1},
  contenu = {
    \cpgeitem{Interféromètre de Michelson éclairé par une source spatialement étendue. Localisation (admise) des franges.}},
  capacites = {
    \cpgeitem{Citer les conditions d’éclairage et d’observation en lame d’air et en coin d’air.}},
  experimentales = {},
  numeriques = {},
}

\cpgeligne{MP/3/4/2}{
  parent = {MP/3/4},
  numero = {3.4.2},
  contenu = {
    \cpgeitem{Lame d’air : franges d’égale inclinaison.}},
  capacites = {
    \cpgeitem{Établir et utiliser l’expression de la différence de marche en fonction de l’épaisseur de la lame d’air équivalente et de l’angle d’incidence des rayons.}},
  experimentales = {
    \cpgeitem{Régler un interféromètre de Michelson pour une observation en lame d’air avec une source étendue à l’aide d’un protocole proposé.\cpgerenvoi{MP/E/2/8}}
    \cpgeitem{Mettre en œuvre un protocole pour accéder au profil spectral d’une raie ou d'un doublet à l’aide d’un interféromètre de Michelson.\cpgerenvoi{MP/E/2/8}}},
  numeriques = {},
}

\cpgeligne{MP/3/4/3}{
  parent = {MP/3/4},
  numero = {3.4.3},
  contenu = {
    \cpgeitem{Coin d’air : franges d’égale épaisseur.}},
  capacites = {
    \cpgeitem{Utiliser l’expression admise de la différence de marche en fonction de l’épaisseur.}},
  experimentales = {
    \cpgeitem{Caractériser la géométrie d’un objet ou l’indice d’un milieu à l’aide d’un interféromètre de Michelson.}
    \cpgeitem{Interpréter qualitativement les observations en lumière blanche.}},
  numeriques = {},
}

\cpgerubrique{MP/4}{
  parent = {MP},
  numero = {4},
  niveau = {1},
  titre = {Électromagnétisme},
  chapeau = {
    \cpgeitem{Le programme d’électromagnétisme de la classe de MP s’inscrit dans le prolongement du thème « Ondes et signaux » du programme de MPSI. Il s’agit pour les étudiants de découvrir les lois locales et intégrales qui gouvernent les champs électrique et magnétique et les phénomènes que ces lois permettent de modéliser, notamment dans le domaine des ondes électromagnétiques.}
    \cpgeitem{L’étude des champs électrostatique et magnétostatique est présentée en deux parties distinctes ; l'enseignant est libre, s’il le souhaite, de procéder à une présentation unifiée de la notion de champ statique. Pour les calculs de champs, l’accent est mis sur les situations à haut degré de symétrie qui permettent l’utilisation efficace des propriétés de flux ou de circulation. Les équations locales des champs statiques sont introduites comme cas particuliers des équations de Maxwell.}
    \cpgeitem{La loi de Biot et Savart, les notions de potentiel vecteur et d’angle solide ne relèvent pas du programme.}
    \cpgeitem{Les relations de passage relatives au champ électromagnétique peuvent être exploitées mais doivent être systématiquement rappelées.}},
}

\cpgerubrique{MP/4/1}{
  parent = {MP/4},
  numero = {4.1},
  niveau = {2},
  titre = {Électrostatique},
  chapeau = {
    \cpgeitem{La partie \cpgecettepartie{} constitue un approfondissement des lois quantitatives qui régissent le champ électrostatique. Les notions abordées sont donc centrées sur l’essentiel : distributions de charges, champ et potentiel. Pour le champ électrostatique et le potentiel, on se limite aux expressions dans le cas de charges ponctuelles.}
    \cpgeitem{L’accent est mis sur les propriétés intégrales du champ et sur le théorème de Gauss pour des situations présentant un haut degré de symétrie ; ce dernier est admis.}
    \cpgeitem{Des capacités sur la lecture des lignes de champ et des surfaces équipotentielles sont développées.}
    \cpgeitem{Le condensateur plan est introduit mais l’étude des conducteurs en équilibre électrostatique ne relève pas du programme.}
    \cpgeitem{Une approche énergétique est conduite dans un cas simple : une charge ponctuelle placée dans un champ électrostatique extérieur.}
    \cpgeitem{Le dipôle est traité, l’accent est mis sur les effets qualitatifs.}
    \cpgeitem{Les analogies avec la gravitation sont centrées sur l’application du théorème de Gauss.}},
}

\cpgeligne{MP/4/1/1}{
  parent = {MP/4/1},
  numero = {4.1.1},
  contenu = {
    \cpgeitem{Loi de Coulomb. Champ électrostatique. Champ électrostatique créé par un ensemble de charges ponctuelles. Principe de superposition.}},
  capacites = {
    \cpgeitem{Exprimer le champ électrostatique créé par une distribution discrète de charges.}
    \cpgeitem{Citer quelques ordres de grandeur de valeurs de champs électrostatiques.}},
  experimentales = {},
  numeriques = {},
}

\cpgeligne{MP/4/1/2}{
  parent = {MP/4/1},
  numero = {4.1.2},
  contenu = {
    \cpgeitem{Distributions continues de charges : volumique, surfacique, linéique.}},
  capacites = {
    \cpgeitem{Choisir un type de distribution continue adaptée à la situation modélisée.}
    \cpgeitem{Relier les densités de charges de deux types de distributions modélisant une même situation.}
    \cpgeitem{Déterminer la charge totale d'une distribution continue dans des situations simples.}},
  experimentales = {},
  numeriques = {},
}

\cpgeligne{MP/4/1/3}{
  parent = {MP/4/1},
  numero = {4.1.3},
  contenu = {
    \cpgeitem{Symétries et invariances du champ électrostatique.}},
  capacites = {
    \cpgeitem{Identifier les plans de symétrie et d'antisymétrie d'une distribution de charges.}
    \cpgeitem{Identifier les invariances d'une distribution de charges.}
    \cpgeitem{Exploiter les symétries et les invariances d'une distribution de charges pour caractériser le champ électrostatique créé.}},
  experimentales = {},
  numeriques = {},
}

\cpgeligne{MP/4/1/4}{
  parent = {MP/4/1},
  numero = {4.1.4},
  contenu = {
    \cpgeitem{Circulation du champ électrostatique.}
    \cpgeitem{Potentiel électrostatique.}
    \cpgeitem{Opérateur gradient.}},
  capacites = {
    \cpgeitem{Relier le champ électrostatique au potentiel. Exprimer le potentiel créé par une distribution discrète de charges.}
    \cpgeitem{Citer l'expression de l'opérateur gradient en coordonnées cartésiennes.}
    \cpgeitem{Déterminer un champ électrostatique à partir du potentiel, l'expression de l'opérateur gradient étant fournie dans le cas des coordonnées sphériques et cylindriques.}
    \cpgeitem{Déterminer une différence de potentiel par circulation du champ électrostatique dans des cas simples.}},
  experimentales = {},
  numeriques = {},
}

\cpgeligne{MP/4/1/5}{
  parent = {MP/4/1},
  numero = {4.1.5},
  contenu = {
    \cpgeitem{Flux du champ électrostatique. Théorème de Gauss.}},
  capacites = {
    \cpgeitem{Identifier les situations pour lesquelles le champ électrostatique peut être calculé à l'aide du théorème de Gauss.}},
  experimentales = {},
  numeriques = {},
}

\cpgeligne{MP/4/1/6}{
  parent = {MP/4/1},
  numero = {4.1.6},
  contenu = {
    \cpgeitem{Systèmes modélisés par une sphère, un cylindre infini ou un plan infini.}},
  capacites = {
    \cpgeitem{Établir les expressions des champs électrostatiques créés en tout point de l'espace par une sphère uniformément chargée en volume, par un cylindre infini uniformément chargé en volume et par un plan infini uniformément chargé en surface.}
    \cpgeitem{Établir et énoncer qu'à l'extérieur d'une distribution à symétrie sphérique, le champ électrostatique créé est le même que celui d'une charge ponctuelle concentrant la charge totale et placée au centre de la distribution.}
    \cpgeitem{Utiliser le théorème de Gauss pour déterminer le champ électrostatique créé par une distribution présentant un haut degré de symétrie.}},
  experimentales = {},
  numeriques = {},
}

\cpgeligne{MP/4/1/7}{
  parent = {MP/4/1},
  numero = {4.1.7},
  contenu = {
    \cpgeitem{Étude du condensateur plan modélisé comme la superposition de deux distributions surfaciques, de charges opposées.}},
  capacites = {
    \cpgeitem{Établir et citer l'expression de la capacité d'un condensateur plan dans le vide.}},
  experimentales = {},
  numeriques = {},
}

\cpgeligne{MP/4/1/8}{
  parent = {MP/4/1},
  numero = {4.1.8},
  contenu = {
    \cpgeitem{Lignes de champ, tubes de champ, surfaces équipotentielles.}},
  capacites = {
    \cpgeitem{Orienter les lignes de champ électrostatique créées par une distribution de charges.}
    \cpgeitem{Représenter les surfaces équipotentielles connaissant les lignes de champ et inversement. Associer les variations de l'intensité du champ électrostatique à la position relative des lignes de champ.}
    \cpgeitem{Vérifier qu'une carte de lignes de champ est compatible avec les symétries et les invariances d'une distribution.}},
  experimentales = {},
  numeriques = {
    \cpgeitem{à l’aide d’un langage de programmation, tracer quelques lignes de champ pour une distribution donnée.}},
}

\cpgeligne{MP/4/1/9}{
  parent = {MP/4/1},
  numero = {4.1.9},
  contenu = {
    \cpgeitem{Énergie potentielle électrostatique d’une charge placée dans un champ électrostatique}
    \cpgeitem{extérieur.}},
  capacites = {
    \cpgeitem{Établir et exploiter l’expression de l'énergie potentielle d'une charge ponctuelle placée dans un champ électrostatique extérieur.}},
  experimentales = {},
  numeriques = {},
}

\cpgeligne{MP/4/1/10}{
  parent = {MP/4/1},
  numero = {4.1.10},
  contenu = {
    \cpgeitem{Notion de dipôle électrostatique, moment dipolaire.}},
  capacites = {
    \cpgeitem{Exprimer le moment dipolaire d'un doublet de charges.}
    \cpgeitem{Évaluer des ordres de grandeur dans le domaine microscopique.}},
  experimentales = {},
  numeriques = {},
}

\cpgeligne{MP/4/1/11}{
  parent = {MP/4/1},
  numero = {4.1.11},
  contenu = {
    \cpgeitem{Champ et potentiel créés par un dipôle électrostatique.}},
  capacites = {
    \cpgeitem{Expliciter l’approximation dipolaire. Représenter l'allure des lignes de champ et des surfaces équipotentielles d'un dipôle électrostatique.}
    \cpgeitem{Établir et exploiter les expressions du champ et du potentiel créés par un doublet de charges dans l’approximation dipolaire.}},
  experimentales = {},
  numeriques = {},
}

\cpgeligne{MP/4/1/12}{
  parent = {MP/4/1},
  numero = {4.1.12},
  contenu = {
    \cpgeitem{Dipôle électrostatique placé dans un champ électrostatique extérieur : actions subies et énergie potentielle d’interaction.}},
  capacites = {
    \cpgeitem{Expliquer qualitativement le comportement d'un dipôle placé dans un champ électrostatique extérieur.}
    \cpgeitem{Établir et exploiter les expressions des actions mécaniques subies par un doublet de charges dans un champ électrostatique extérieur uniforme.}
    \cpgeitem{Exploiter l’expression fournie de la force subie par un dipôle placé dans un champ électrostatique extérieur non uniforme.}
    \cpgeitem{Citer et exploiter l’expression de l’énergie potentielle d’interaction.}},
  experimentales = {},
  numeriques = {},
}

\cpgeligne{MP/4/1/13}{
  parent = {MP/4/1},
  numero = {4.1.13},
  contenu = {
    \cpgeitem{Analogies avec la gravitation.}},
  capacites = {
    \cpgeitem{Utiliser le théorème de Gauss de la gravitation.}},
  experimentales = {},
  numeriques = {},
}

\cpgerubrique{MP/4/2}{
  parent = {MP/4},
  numero = {4.2},
  niveau = {2},
  titre = {Magnétostatique},
  chapeau = {
    \cpgeitem{La partie \cpgecettepartie{} s’appuie sur les différents aspects qualitatifs et quantitatifs vus en première année de la classe de MPSI. Les étudiants sont donc déjà familiarisés avec le concept de champ magnétostatique. La loi de Biot et Savart n’est pas introduite ; l’utilisation de celle-ci pour calculer un champ magnétostatique est donc exclue.}
    \cpgeitem{Les distributions de courants surfaciques ne sont pas introduites à ce niveau du programme, elles le sont uniquement à l’occasion de la réflexion d’une onde électromagnétique sur un métal parfait.}
    \cpgeitem{L’étude des propriétés des dipôles magnétiques, déjà abordée en classe de MPSI est ici complétée notamment en ce qui concerne les actions subies par un dipôle magnétique placé dans un champ magnétostatique extérieur. On peut, sur ce thème, souligner les analogies avec l’électrostatique.}},
}

\cpgeligne{MP/4/2/1}{
  parent = {MP/4/2},
  numero = {4.2.1},
  contenu = {
    \cpgeitem{Vecteur densité de courant volumique. Intensité du courant. Distributions de courant volumique et linéique.}},
  capacites = {
    \cpgeitem{Relier l’intensité du courant et le flux du vecteur densité de courant volumique.}},
  experimentales = {},
  numeriques = {},
}

\cpgeligne{MP/4/2/2}{
  parent = {MP/4/2},
  numero = {4.2.2},
  contenu = {
    \cpgeitem{Symétries et invariances des distributions de courant.}},
  capacites = {
    \cpgeitem{Exploiter les propriétés de symétrie et d’invariance des sources pour prévoir des propriétés du champ créé.}},
  experimentales = {},
  numeriques = {},
}

\cpgeligne{MP/4/2/3}{
  parent = {MP/4/2},
  numero = {4.2.3},
  contenu = {
    \cpgeitem{Propriétés de flux et de circulation. Théorème d’Ampère.}},
  capacites = {
    \cpgeitem{Identifier les situations pour lesquelles le champ magnétostatique peut être calculé à l'aide du théorème d'Ampère.}
    \cpgeitem{Choisir un contour, une surface et les orienter pour appliquer le théorème d’Ampère en vue de déterminer l’expression d’un champ magnétique. Utiliser une méthode de superposition.}
    \cpgeitem{Citer quelques ordres de grandeur de valeurs de champs magnétostatiques.}},
  experimentales = {},
  numeriques = {},
}

\cpgeligne{MP/4/2/4}{
  parent = {MP/4/2},
  numero = {4.2.4},
  contenu = {
    \cpgeitem{Modèles du fil rectiligne infini de section non nulle et du solénoïde infini.}},
  capacites = {
    \cpgeitem{Établir les expressions des champs magnétostatiques créés en tout point de l'espace par un fil rectiligne infini de section non nulle, parcouru par des courants uniformément répartis en volume, par un solénoïde infini en admettant que le champ est nul à l’extérieur.}},
  experimentales = {},
  numeriques = {},
}

\cpgeligne{MP/4/2/5}{
  parent = {MP/4/2},
  numero = {4.2.5},
  contenu = {
    \cpgeitem{Lignes de champ, tubes de champ.}},
  capacites = {
    \cpgeitem{Orienter les lignes de champ magnétostatique créées par une distribution de courants.}
    \cpgeitem{Associer les variations de l'intensité du champ magnétostatique à la position relative des lignes de champ.}
    \cpgeitem{Vérifier qu'une carte de lignes de champ est compatible avec les symétries et les invariances d'une distribution}},
  experimentales = {},
  numeriques = {},
}

\cpgeligne{MP/4/2/6}{
  parent = {MP/4/2},
  numero = {4.2.6},
  contenu = {
    \cpgeitem{Notion de dipôle magnétique. Moment magnétique.}},
  capacites = {
    \cpgeitem{Exprimer le moment magnétique d’une boucle de courant plane.}
    \cpgeitem{Évaluer des ordres de grandeur dans les domaines macroscopique et microscopique.}},
  experimentales = {},
  numeriques = {},
}

\cpgeligne{MP/4/2/7}{
  parent = {MP/4/2},
  numero = {4.2.7},
  contenu = {
    \cpgeitem{Champ créé par un dipôle magnétique.}},
  capacites = {
    \cpgeitem{Expliciter l’approximation dipolaire. Représenter l'allure des lignes de champ d'un dipôle magnétique.}
    \cpgeitem{Exploiter l’expression fournie du champ créé par un dipôle magnétique.}},
  experimentales = {},
  numeriques = {},
}

\cpgeligne{MP/4/2/8}{
  parent = {MP/4/2},
  numero = {4.2.8},
  contenu = {
    \cpgeitem{Dipôle magnétique placé dans un champ magnétostatique extérieur : actions subies et énergie potentielle d’interaction.}},
  capacites = {
    \cpgeitem{Expliquer qualitativement le comportement d'un dipôle passif placé dans un champ magnétostatique extérieur.}
    \cpgeitem{Exploiter les expressions fournies des actions mécaniques subies par un dipôle magnétique dans un champ magnétostatique extérieur uniforme.}
    \cpgeitem{Exploiter l’expression fournie de la force subie par un dipôle magnétique dans un champ magnétostatique extérieur non uniforme.}
    \cpgeitem{Citer et exploiter l’expression de l’énergie potentielle d’interaction.}},
  experimentales = {},
  numeriques = {},
}

\cpgerubrique{MP/4/3}{
  parent = {MP/4},
  numero = {4.3},
  niveau = {2},
  titre = {Équations de Maxwell},
  chapeau = {
    \cpgeitem{Dans la partie \cpgecettepartie{} une vision unifiée des lois de l’électromagnétisme est présentée. Elle conduit à une première approche quantitative du phénomène de propagation et permet également de revenir sur les lois de l’induction étudiées en première année de MPSI. Les lois locales de l’électrostatique relatives au potentiel constituent un support pertinent pour procéder à une approche numérique de la résolution d’une équation aux dérivées partielles.}},
}

\cpgeligne{MP/4/3/1}{
  parent = {MP/4/3},
  numero = {4.3.1},
  contenu = {
    \cpgeitem{Principe de la conservation de la charge : formulation locale.}},
  capacites = {
    \cpgeitem{Établir l’équation locale de la conservation de la charge en coordonnées cartésiennes dans le cas à une dimension.}},
  experimentales = {},
  numeriques = {},
}

\cpgeligne{MP/4/3/2}{
  parent = {MP/4/3},
  numero = {4.3.2},
  contenu = {
    \cpgeitem{Équations de Maxwell : formulations locale et intégrale.}},
  capacites = {
    \cpgeitem{Associer l’équation de Maxwell-Faraday à la loi de Faraday.}
    \cpgeitem{Citer, utiliser et interpréter les équations de Maxwell sous forme intégrale.}
    \cpgeitem{Associer qualitativement le couplage spatio-temporel entre champ électrique et champ magnétique au phénomène de propagation. Vérifier la cohérence des équations de Maxwell avec l’équation locale de la conservation de la charge.}},
  experimentales = {},
  numeriques = {},
}

\cpgeligne{MP/4/3/3}{
  parent = {MP/4/3},
  numero = {4.3.3},
  contenu = {
    \cpgeitem{Équations de propagation des champs dans une région vide de charges et de courants.}},
  capacites = {
    \cpgeitem{Établir les équations de propagation à partir des équations de Maxwell.}},
  experimentales = {},
  numeriques = {},
}

\cpgeligne{MP/4/3/4}{
  parent = {MP/4/3},
  numero = {4.3.4},
  contenu = {
    \cpgeitem{Cas des champs statiques : équations locales.}},
  capacites = {
    \cpgeitem{Établir les lois locales des champs statiques à partir des équations de Maxwell.}},
  experimentales = {},
  numeriques = {},
}

\cpgeligne{MP/4/3/5}{
  parent = {MP/4/3},
  numero = {4.3.5},
  contenu = {
    \cpgeitem{Équation de Poisson et équation de Laplace de l’électrostatique.}},
  capacites = {
    \cpgeitem{Établir les équations de Poisson et de Laplace de l’électrostatique.}
    \cpgeitem{Exprimer par analogie les équations de Poisson et de Laplace dans le cas de la gravitation.}},
  experimentales = {},
  numeriques = {},
}

\cpgerubrique{MP/4/4}{
  parent = {MP/4},
  numero = {4.4},
  niveau = {2},
  titre = {Énergie du champ électromagnétique},
  chapeau = {
    \cpgeitem{Dans la partie \cpgecettepartie{}, aucun modèle relatif à la loi d’Ohm locale n’est exigible ; l’accent est mis sur les échanges d’énergie entre la matière et le champ électromagnétique, sur l’utilisation du flux du vecteur de Poynting pour évaluer une puissance rayonnée à travers une surface et sur les bilans d’énergie et de puissance.}},
}

\cpgeligne{MP/4/4/1}{
  parent = {MP/4/4},
  numero = {4.4.1},
  contenu = {
    \cpgeitem{Force électromagnétique volumique. Puissance volumique cédée par le champ électromagnétique aux porteurs de charge.}},
  capacites = {
    \cpgeitem{Établir et utiliser l’expression de la puissance volumique cédée par le champ électromagnétique aux porteurs de charge.}},
  experimentales = {},
  numeriques = {},
}

\cpgeligne{MP/4/4/2}{
  parent = {MP/4/4},
  numero = {4.4.2},
  contenu = {
    \cpgeitem{Loi d’Ohm locale ; puissance volumique dissipée par effet Joule.}},
  capacites = {
    \cpgeitem{Analyser les aspects énergétiques dans le cas particulier d’un milieu ohmique.}},
  experimentales = {},
  numeriques = {},
}

\cpgeligne{MP/4/4/3}{
  parent = {MP/4/4},
  numero = {4.4.3},
  contenu = {
    \cpgeitem{Énergie électromagnétique volumique.}
    \cpgeitem{Vecteur de Poynting. Bilan d’énergie.}},
  capacites = {
    \cpgeitem{Citer des ordres de grandeur de flux énergétiques moyens (flux solaire, laser…).}
    \cpgeitem{Utiliser le flux du vecteur de Poynting à travers une surface orientée pour évaluer la puissance rayonnée.}
    \cpgeitem{Effectuer un bilan d’énergie sous forme locale et intégrale.}
    \cpgeitem{Interpréter chaque terme de l’équation locale de Poynting, celle-ci étant fournie.}},
  experimentales = {},
  numeriques = {},
}

\cpgerubrique{MP/4/5}{
  parent = {MP/4},
  numero = {4.5},
  niveau = {2},
  titre = {Propagation et rayonnement},
  chapeau = {
    \cpgeitem{La partie \cpgecettepartie{} est l’occasion d’illustrer l’efficacité du formalisme local des équations de Maxwell en insistant sur les aspects qualitatifs et sur la variété des applications qui en découlent.}
    \cpgeitem{Si le modèle de l’onde plane est présenté dans le cadre de l’espace vide, les études des ondes électromagnétiques dans un plasma ainsi que dans un milieu ohmique permettent d’illustrer l’importance des couplages entre les champs, les charges et les courants. Elles sont également l’occasion d’enrichir les compétences des étudiants sur les phénomènes de propagation en   abordant, par exemple, l’effet de peau, le phénomène de dispersion, les notions de vitesse de groupe et de phase, de fréquence de coupure ou encore d’onde évanescente.}
    \cpgeitem{La réflexion d’une onde électromagnétique sur un métal parfait permet d’aborder la notion d’onde stationnaire. L’importance des conditions aux limites imposées sur la quantification des solutions est soulignée. La notion de densité de courant surfacique est introduite mais le calcul de l’intensité à travers un segment ne relève pas du programme.}
    \cpgeitem{L’étude du rayonnement dipolaire repose sur l’analyse et l’exploitation des expressions des champs, qui sont admises. Elle est l’occasion d’étudier une modélisation du phénomène de diffusion d’une onde électromagnétique par un atome et d’en analyser les conséquences.}},
}

\cpgeligne{MP/4/5/1}{
  parent = {MP/4/5},
  numero = {4.5.1},
  contenu = {
    \cpgeitem{Onde plane dans l’espace vide de charge et de courant ; onde plane progressive et aspects énergétiques.}},
  capacites = {
    \cpgeitem{Citer les solutions de l’équation de d’Alembert à une dimension.}
    \cpgeitem{Décrire la structure d’une onde plane et d’une onde plane progressive dans l’espace vide de charge et de courant.}},
  experimentales = {},
  numeriques = {},
}

\cpgeligne{MP/4/5/2}{
  parent = {MP/4/5},
  numero = {4.5.2},
  contenu = {
    \cpgeitem{Onde plane progressive monochromatique. Relation de dispersion.}},
  capacites = {
    \cpgeitem{Expliquer le caractère idéal du modèle de l’onde plane monochromatique.}
    \cpgeitem{Déterminer la relation de dispersion. Citer les domaines du spectre des ondes électromagnétiques et leur associer des applications.}
    \cpgeitem{Exprimer le vecteur de Poynting et l’énergie électromagnétique volumique associés à une onde plane progressive monochromatique.}
    \cpgeitem{Effectuer une étude énergétique dans le cas d’une onde plane progressive monochromatique.}},
  experimentales = {},
  numeriques = {},
}

\cpgeligne{MP/4/5/3}{
  parent = {MP/4/5},
  numero = {4.5.3},
  contenu = {
    \cpgeitem{Onde plane progressive monochromatique polarisée rectilignement ou circulairement.}},
  capacites = {
    \cpgeitem{Reconnaître une onde polarisée rectilignement ou circulairement.}},
  experimentales = {
    \cpgeitem{Utiliser des polariseurs et étudier quantitativement la loi de Malus.\cpgerenvoi{MP/E/2/6}}},
  numeriques = {},
}

\cpgeligne{MP/4/5/4}{
  parent = {MP/4/5},
  numero = {4.5.4},
  contenu = {
    \cpgeitem{Onde plane transverse électrique monochromatique dans un plasma dilué.}
    \cpgeitem{Conductivité complexe du milieu.}
    \cpgeitem{Pulsation de coupure.}
    \cpgeitem{Ondes évanescentes.}},
  capacites = {
    \cpgeitem{Exprimer la conductivité complexe du milieu et établir la relation de dispersion.}
    \cpgeitem{Décrire le phénomène de dispersion.}
    \cpgeitem{Relier la fréquence de coupure aux caractéristiques du plasma et citer son ordre de grandeur dans le cas de l’ionosphère.}
    \cpgeitem{Distinguer qualitativement les ondes évanescentes et les ondes progressives du point de vue du transport de l’énergie.}},
  experimentales = {},
  numeriques = {},
}

\cpgeligne{MP/4/5/5}{
  parent = {MP/4/5},
  numero = {4.5.5},
  contenu = {
    \cpgeitem{Vitesse de phase, vitesse de groupe.}
    \cpgeitem{Propagation d’un paquet d’ondes dans un milieu linéaire faiblement dispersif.}},
  capacites = {
    \cpgeitem{Calculer la vitesse de groupe à partir de la relation de dispersion.}
    \cpgeitem{Associer la vitesse de groupe à la propagation de l’enveloppe du paquet d’ondes.}},
  experimentales = {},
  numeriques = {},
}

\cpgeligne{MP/4/5/6}{
  parent = {MP/4/5},
  numero = {4.5.6},
  contenu = {
    \cpgeitem{Propagation d’une onde électromagnétique dans un milieu ohmique en régime lentement}
    \cpgeitem{variable. Effet de peau.}},
  capacites = {
    \cpgeitem{Établir et interpréter l’expression de la longueur caractéristique d’atténuation de l’onde électromagnétique dans un milieu ohmique.}},
  experimentales = {},
  numeriques = {},
}

\cpgeligne{MP/4/5/7}{
  parent = {MP/4/5},
  numero = {4.5.7},
  contenu = {
    \cpgeitem{Réflexion sous incidence normale d’une onde plane, progressive et monochromatique polarisée rectilignement sur un plan conducteur parfait. Onde stationnaire.}},
  capacites = {
    \cpgeitem{Établir l’expression de l’onde réfléchie en exploitant les relations de passage fournies. Interpréter qualitativement la présence de courants localisés en surface.}
    \cpgeitem{Reconnaître et caractériser une onde stationnaire.}},
  experimentales = {},
  numeriques = {},
}

\cpgeligne{MP/4/5/8}{
  parent = {MP/4/5},
  numero = {4.5.8},
  contenu = {
    \cpgeitem{Applications aux cavités à une dimension.}
    \cpgeitem{Mode d’onde stationnaire.}},
  capacites = {
    \cpgeitem{Établir la condition de quantification des solutions.}},
  experimentales = {
    \cpgeitem{Mettre en œuvre un dispositif permettant d’étudier une onde électromagnétique, dans le domaine des ondes centimétriques.\cpgerenvoi{MP/E/2/5}}},
  numeriques = {},
}

\cpgeligne{MP/4/5/9}{
  parent = {MP/4/5},
  numero = {4.5.9},
  contenu = {
    \cpgeitem{Champ électromagnétique rayonné par un dipôle oscillant dans la zone de rayonnement.}
    \cpgeitem{Puissance rayonnée.}},
  capacites = {
    \cpgeitem{Justifier l’intérêt du modèle du dipôle oscillant et citer des exemples dans différents domaines.}
    \cpgeitem{Formuler et commenter les approximations reliant les trois échelles de longueur pertinentes.}
    \cpgeitem{Analyser la structure du champ électromagnétique rayonné, les expressions des champs étant fournies, en utilisant des arguments généraux : symétrie, conservation de l’énergie et analyse dimensionnelle.}
    \cpgeitem{Effectuer un bilan énergétique, les expressions des champs étant fournies.}
    \cpgeitem{Représenter l’indicatrice de rayonnement.}},
  experimentales = {
    \cpgeitem{Détecter une onde électromagnétique rayonnée.}},
  numeriques = {},
}

\cpgeligne{MP/4/5/10}{
  parent = {MP/4/5},
  numero = {4.5.10},
  contenu = {
    \cpgeitem{Diffusion d’une onde électromagnétique polarisée rectilignement par une molécule dans cadre du modèle de la charge élastiquement liée.}
    \cpgeitem{Structure de l’onde diffusée.}
    \cpgeitem{Puissance diffusée en fonction de la fréquence. Résonance. Domaine de}
    \cpgeitem{Rayleigh.}},
  capacites = {
    \cpgeitem{Déterminer les caractéristiques du dipôle induit en régime établi, par l’action de l’onde incidente sur la molécule.}
    \cpgeitem{Identifier les domaines de résonances et de Rayleigh.}
    \cpgeitem{Citer des illustrations de la diffusion d’une onde électromagnétique par un milieu.}},
  experimentales = {},
  numeriques = {},
}

\cpgerubrique{MP/5}{
  parent = {MP},
  numero = {5},
  niveau = {1},
  titre = {Thermodynamique},
  chapeau = {
    \cpgeitem{Le programme de thermodynamique de la classe de MP s’inscrit dans le prolongement de celui de MPSI. Les principes de la thermodynamique sont désormais écrits sous forme infinitésimale $\cpgedd U+\cpgedd E = \delta W+\delta Q$ et $\cpgedd S = \delta S_e +\delta S_c$. Le premier principe écrit sous forme infinitésimale est réinvesti dans l’étude des transferts thermiques.}},
}

\cpgerubrique{MP/5/1}{
  parent = {MP/5},
  numero = {5.1},
  niveau = {2},
  titre = {Principes de la thermodynamique},
  chapeau = {
    \cpgeitem{La partie intitulée \cpgecettepartie{} est principalement consacrée à l’étude des systèmes ouverts limitée au cas d’un écoulement unidimensionnel.}},
}

\cpgeligne{MP/5/1/1}{
  parent = {MP/5/1},
  numero = {5.1.1},
  contenu = {
    \cpgeitem{Formulation des principes de la thermodynamique pour une transformation élémentaire.}},
  capacites = {
    \cpgeitem{Utiliser avec rigueur les notations $\cpgedd$ et $\delta$ en leur attachant une signification.}},
  experimentales = {},
  numeriques = {},
}

\cpgeligne{MP/5/1/2}{
  parent = {MP/5/1},
  numero = {5.1.2},
  contenu = {
    \cpgeitem{Premier et deuxième principes de la thermodynamique pour un système ouvert en régime stationnaire, dans le seul cas d’un écoulement unidimensionnel dans la section}
    \cpgeitem{d’entrée et la section de sortie.}},
  capacites = {
    \cpgeitem{Établir les relations $\Delta h+\Delta e = w_u+q$ et $\Delta s = s_e+s_c$ et les utiliser pour étudier des machines thermiques réelles à l’aide du diagramme $(P,h)$.}},
  experimentales = {},
  numeriques = {},
}

\cpgerubrique{MP/5/2}{
  parent = {MP/5},
  numero = {5.2},
  niveau = {2},
  titre = {Transferts thermiques},
  chapeau = {
    \cpgeitem{Dans la partie \cpgecettepartie{}, l’établissement de l’équation de la diffusion thermique est limitée au cas des solides ; il est possible d’utiliser les résultats établis dans d’autres situations, notamment dans le cas de l’étude des fluides, en affirmant la généralisation des équations obtenues dans le cas des solides. Les mises en équations locales sont faites exclusivement sur des géométries où une seule variable d’espace intervient. On admet ensuite les formes générales des équations en utilisant les opérateurs d’analyse vectorielle. Enfin, aucune connaissance spécifique sur les solutions d’une équation de diffusion ne figure au programme.}
    \cpgeitem{La loi de Newton à l’interface entre un solide et un fluide est introduite.}},
}

\cpgeligne{MP/5/2/1}{
  parent = {MP/5/2},
  numero = {5.2.1},
  contenu = {
    \cpgeitem{Conduction, convection et rayonnement.}},
  capacites = {
    \cpgeitem{Identifier un mode de transfert thermique.}},
  experimentales = {
    \cpgeitem{Mettre en œuvre un dispositif expérimental utilisant une caméra thermique ou un capteur dans le domaine des infrarouges.}},
  numeriques = {},
}

\cpgeligne{MP/5/2/2}{
  parent = {MP/5/2},
  numero = {5.2.2},
  contenu = {
    \cpgeitem{Flux thermique.}
    \cpgeitem{Vecteur densité de flux thermique.}},
  capacites = {
    \cpgeitem{Calculer un flux thermique à travers une surface orientée et interpréter son signe.}},
  experimentales = {},
  numeriques = {},
}

\cpgeligne{MP/5/2/3}{
  parent = {MP/5/2},
  numero = {5.2.3},
  contenu = {
    \cpgeitem{Premier principe de la thermodynamique.}},
  capacites = {
    \cpgeitem{Effectuer un bilan local d’énergie interne pour un solide dans le cas d’une situation à une variable d’espace en géométrie cartésienne, cylindrique}
    \cpgeitem{ou sphérique.}},
  experimentales = {},
  numeriques = {},
}

\cpgeligne{MP/5/2/4}{
  parent = {MP/5/2},
  numero = {5.2.4},
  contenu = {
    \cpgeitem{Loi de Fourier.}},
  capacites = {
    \cpgeitem{Interpréter et utiliser la loi de Fourier.}
    \cpgeitem{Citer quelques ordres de grandeur de conductivité thermique dans les conditions usuelles : air, eau, verre, acier.}},
  experimentales = {
    \cpgeitem{Mesurer la conductivité thermique d’un matériau.}},
  numeriques = {},
}

\cpgeligne{MP/5/2/5}{
  parent = {MP/5/2},
  numero = {5.2.5},
  contenu = {
    \cpgeitem{Équation de la diffusion thermique.}},
  capacites = {
    \cpgeitem{Établir l’équation de la diffusion thermique sans terme de source au sein d’un solide dans le cas d’une situation à une variable d’espace en géométrie cartésienne, cylindrique ou sphérique. Utiliser une généralisation de l’équation de la diffusion en présence d’un terme de source.}
    \cpgeitem{Utiliser une généralisation en géométrie quelconque en utilisant l’opérateur Laplacien et son expression fournie.}
    \cpgeitem{Analyser une équation de diffusion thermique en ordre de grandeur pour relier des échelles caractéristiques spatiale et temporelle.}},
  experimentales = {},
  numeriques = {
    \cpgeitem{à l’aide d’un langage de programmation, résoudre l’équation de la diffusion thermique à une dimension par une méthode des différences finies dérivée de la méthode d’Euler explicite de résolution des équations différentielles ordinaires.}},
}

\cpgeligne{MP/5/2/6}{
  parent = {MP/5/2},
  numero = {5.2.6},
  contenu = {
    \cpgeitem{Régime stationnaire. Résistance thermique.}},
  capacites = {
    \cpgeitem{Définir la notion de résistance thermique par analogie avec l’électrocinétique.}
    \cpgeitem{Déterminer l’expression de la résistance thermique d’un solide dans le cas d’un problème unidimensionnel en géométrie cartésienne.}
    \cpgeitem{Exploiter les lois d’association de résistances thermiques.}},
  experimentales = {},
  numeriques = {},
}

\cpgeligne{MP/5/2/7}{
  parent = {MP/5/2},
  numero = {5.2.7},
  contenu = {
    \cpgeitem{Coefficient de transfert thermique de surface ; loi de Newton.}},
  capacites = {
    \cpgeitem{Utiliser la loi de Newton comme condition aux limites à une interface solide-fluide.}},
  experimentales = {},
  numeriques = {},
}

\cpgerubrique{MP/6}{
  parent = {MP},
  numero = {6},
  niveau = {1},
  titre = {Physique quantique},
  chapeau = {
    \cpgeitem{Cette partie s’inscrit dans le prolongement du programme de la classe de MPSI. Il s’agit cependant de dépasser l’approche descriptive et qualitative et de donner aux étudiants leurs premiers outils quantitatifs d’analyse. Le cœur de cet enseignement est construit sur la mécanique ondulatoire de Schrödinger et propose des résolutions d’exemples simples mais fondamentaux pour la bonne compréhension de problèmes plus complexes : particule dans une marche de potentiel et effet tunnel, particule dans un puits de potentiel infini et quantification de l’énergie d’une particule confinée. On se limite à l’introduction heuristique de la dualité onde-particule et de la densité de courant de probabilité pour une particule libre sans développer la notion de paquet d’ondes.}
    \cpgeitem{L’accent doit être mis sur l’interprétation et l’exploitation des résultats et non pas sur les calculs, non exigibles pour l’exemple plus délicat de la barrière de potentiel. Le professeur peut au contraire, s’il le souhaite, proposer des analyses de graphes, des exploitations de formules analytiques fournies, des estimations numériques, des simulations… afin d’aborder des modélisations plus réalistes.}},
}

\cpgerubrique{MP/6/1}{
  parent = {MP/6},
  numero = {6.1},
  niveau = {2},
  titre = {Fonction d’onde et équation de Schrödinger},
  chapeau = {},
}

\cpgeligne{MP/6/1/1}{
  parent = {MP/6/1},
  numero = {6.1.1},
  contenu = {
    \cpgeitem{Fonction d’onde $\psi$ d’une particule sans spin et densité de probabilité de présence.}},
  capacites = {
    \cpgeitem{Interpréter en termes de probabilité l’amplitude d’une onde associée à une particule.}},
  experimentales = {},
  numeriques = {},
}

\cpgeligne{MP/6/1/2}{
  parent = {MP/6/1},
  numero = {6.1.2},
  contenu = {
    \cpgeitem{Équation de Schrödinger à une dimension dans un potentiel $V(x)$.}},
  capacites = {
    \cpgeitem{Utiliser le caractère linéaire de l’équation (principe de superposition).}},
  experimentales = {},
  numeriques = {},
}

\cpgeligne{MP/6/1/3}{
  parent = {MP/6/1},
  numero = {6.1.3},
  contenu = {
    \cpgeitem{États stationnaires de l’équation de Schrödinger.}},
  capacites = {
    \cpgeitem{Procéder à la séparation des variables temps et espace.}
    \cpgeitem{Distinguer l’onde associée à un état stationnaire en mécanique quantique d’une onde stationnaire au sens usuel de la physique des ondes.}
    \cpgeitem{Relier l’énergie de la particule à l’évolution temporelle de sa fonction d’onde et faire le lien avec la relation de Planck-Einstein.}
    \cpgeitem{Identifier le terme associé à l’énergie cinétique.}},
  experimentales = {},
  numeriques = {},
}

\cpgerubrique{MP/6/2}{
  parent = {MP/6},
  numero = {6.2},
  niveau = {2},
  titre = {Particule libre},
  chapeau = {},
}

\cpgeligne{MP/6/2/1}{
  parent = {MP/6/2},
  numero = {6.2.1},
  contenu = {
    \cpgeitem{Fonction d’onde d’une particule libre non localisée.}},
  capacites = {
    \cpgeitem{Établir les solutions.}
    \cpgeitem{Interpréter la difficulté de normalisation de cette fonction d’onde.}},
  experimentales = {},
  numeriques = {},
}

\cpgeligne{MP/6/2/2}{
  parent = {MP/6/2},
  numero = {6.2.2},
  contenu = {
    \cpgeitem{Relation de de Broglie.}},
  capacites = {
    \cpgeitem{Relier l’énergie de la particule et le vecteur d’onde de l’onde plane associée.}},
  experimentales = {},
  numeriques = {},
}

\cpgeligne{MP/6/2/3}{
  parent = {MP/6/2},
  numero = {6.2.3},
  contenu = {
    \cpgeitem{Inégalité d’Heisenberg spatiale et paquet d’ondes.}},
  capacites = {
    \cpgeitem{Expliquer, en s’appuyant sur l’inégalité d’Heisenberg spatiale, que la localisation de la particule peut s’obtenir par superposition d’ondes planes.}},
  experimentales = {},
  numeriques = {},
}

\cpgeligne{MP/6/2/4}{
  parent = {MP/6/2},
  numero = {6.2.4},
  contenu = {
    \cpgeitem{Densité de courant de probabilité associée à une particule libre.}},
  capacites = {
    \cpgeitem{Utiliser l’expression admise du courant de probabilité associé à une particule libre ; l’interpréter comme un produit densité$\times$vitesse.}},
  experimentales = {},
  numeriques = {},
}

\cpgerubrique{MP/6/3}{
  parent = {MP/6},
  numero = {6.3},
  niveau = {2},
  titre = {États stationnaires d’une particule dans des potentiels constants par morceaux},
  chapeau = {},
}

\cpgeligne{MP/6/3/1}{
  parent = {MP/6/3},
  numero = {6.3.1},
  contenu = {
    \cpgeitem{États stationnaires d’une particule dans le cas d’une marche de potentiel.}
    \cpgeitem{Cas $E>V$ : probabilité de transmission et de réflexion.}
    \cpgeitem{Cas $E<V$ : évanescence.}},
  capacites = {
    \cpgeitem{Citer des exemples physiques illustrant cette problématique.}
    \cpgeitem{Exploiter les conditions de continuité (admises) relatives à la fonction d’onde.}
    \cpgeitem{Établir la solution dans le cas d’une particule incidente sur une marche de potentiel.}
    \cpgeitem{Expliquer les différences de comportement par rapport à une particule classique}
    \cpgeitem{Déterminer les coefficients de transmission et de réflexion en utilisant les courants de probabilités. Reconnaître l’existence d’une onde évanescente et la caractériser.}},
  experimentales = {},
  numeriques = {},
}

\cpgeligne{MP/6/3/2}{
  parent = {MP/6/3},
  numero = {6.3.2},
  contenu = {
    \cpgeitem{Barrière de potentiel et effet tunnel.}},
  capacites = {
    \cpgeitem{Décrire qualitativement l’influence de la hauteur ou de largeur de la barrière de potentiel sur le coefficient de transmission.}
    \cpgeitem{Exploiter un coefficient de transmission fourni. Citer des applications.}},
  experimentales = {},
  numeriques = {},
}

\cpgeligne{MP/6/3/3}{
  parent = {MP/6/3},
  numero = {6.3.3},
  contenu = {
    \cpgeitem{États stationnaires d’une particule dans un puits de potentiel infini.}},
  capacites = {
    \cpgeitem{Établir les solutions et les niveaux d’énergie de la particule confinée.}
    \cpgeitem{Identifier des analogies avec d’autres domaines de la physique.}},
  experimentales = {},
  numeriques = {},
}

\cpgeligne{MP/6/3/4}{
  parent = {MP/6/3},
  numero = {6.3.4},
  contenu = {
    \cpgeitem{Énergie de confinement.}},
  capacites = {
    \cpgeitem{Estimer l’énergie d’une particule confinée dans son état fondamental pour un puits non rectangulaire.}
    \cpgeitem{Associer l’analyse à l’inégalité d’Heisenberg.}},
  experimentales = {},
  numeriques = {},
}

\cpgerubrique{MP/6/4}{
  parent = {MP/6},
  numero = {6.4},
  niveau = {2},
  titre = {États non stationnaires d’une particule dans un puits de potentiel infini},
  chapeau = {},
}

\cpgeligne{MP/6/4/1}{
  parent = {MP/6/4},
  numero = {6.4.1},
  contenu = {
    \cpgeitem{Combinaison linéaire d’états stationnaires.}},
  capacites = {
    \cpgeitem{Expliquer qu’une superposition de deux états stationnaires engendre une évolution au cours du temps de l’état de la particule.}
    \cpgeitem{Établir l’expression de la densité de probabilité de présence de la particule dans le cas d’une superposition de deux états stationnaires ; interpréter le résultat.}},
  experimentales = {},
  numeriques = {},
}

\cpgerubrique{MP/7}{
  parent = {MP},
  numero = {7},
  niveau = {1},
  titre = {Éléments de thermodynamique statistique},
  chapeau = {
    \cpgeitem{L’objectif de cette partie est de relier certaines propriétés macroscopiques d’un système constitué d’un grand nombre de particules avec celles des constituants microscopiques. Le facteur de Boltzmann est introduit de manière inductive à partir du modèle d’atmosphère isotherme. L’étude des systèmes à spectre discret d’énergies est l’occasion de montrer, qu’à température donnée l’énergie fluctue et que les fluctuations relatives diminuent avec la taille du système. L’étude des systèmes à deux niveaux, conduite de manière plus exhaustive, permet une analyse plus fine des phénomènes.}
    \cpgeitem{L’étude de l’énergie d’un ensemble de particules libres dans une boîte unidimensionnelle à une température donnée établit un lien entre la physique quantique et les propriétés macroscopiques de la matière.}
    \cpgeitem{Le théorème d’équipartition de l’énergie est l’occasion de procéder à une évaluation des capacités thermiques des gaz et des solides.}},
}

\cpgerubrique{MP/7/1}{
  parent = {MP/7},
  numero = {7.1},
  niveau = {2},
  titre = {Facteur de Boltzmann},
  chapeau = {},
}

\cpgeligne{MP/7/1/1}{
  parent = {MP/7/1},
  numero = {7.1.1},
  contenu = {
    \cpgeitem{Modèle de l’atmosphère isotherme.}},
  capacites = {
    \cpgeitem{Établir la variation de la pression avec l’altitude dans l’hypothèse d’une atmosphère isotherme.}},
  experimentales = {},
  numeriques = {},
}

\cpgeligne{MP/7/1/2}{
  parent = {MP/7/1},
  numero = {7.1.2},
  contenu = {
    \cpgeitem{Poids de Boltzmann d’une particule indépendante à l'équilibre avec un thermostat.}},
  capacites = {
    \cpgeitem{Interpréter la loi du nivellement barométrique avec le poids de Boltzmann.}
    \cpgeitem{Identifier un facteur de Boltzmann.}
    \cpgeitem{Comparer $kT$ à des écarts d’énergie et estimer les conséquences d'une variation de température.}},
  experimentales = {},
  numeriques = {},
}

\cpgerubrique{MP/7/2}{
  parent = {MP/7},
  numero = {7.2},
  niveau = {2},
  titre = {Système à spectres discrets d’énergie},
  chapeau = {},
}

\cpgeligne{MP/7/2/1}{
  parent = {MP/7/2},
  numero = {7.2.1},
  contenu = {
    \cpgeitem{Probabilité d'occupation d'un état d'énergie non dégénéré par une particule indépendante.}},
  capacites = {
    \cpgeitem{Exprimer la probabilité d'occupation d’un état d’énergie en utilisant la condition de normalisation.}
    \cpgeitem{Exploiter un rapport de probabilités entre deux états.}},
  experimentales = {},
  numeriques = {},
}

\cpgeligne{MP/7/2/2}{
  parent = {MP/7/2},
  numero = {7.2.2},
  contenu = {
    \cpgeitem{Énergie moyenne et écart quadratique moyen.}},
  capacites = {
    \cpgeitem{Exprimer sous forme d'une somme sur ses états l'énergie moyenne et l’écart-quadratique}
    \cpgeitem{énergétique d'un système.}},
  experimentales = {},
  numeriques = {},
}

\cpgeligne{MP/7/2/3}{
  parent = {MP/7/2},
  numero = {7.2.3},
  contenu = {
    \cpgeitem{Cas d'un système à $N$ particules indépendantes.}},
  capacites = {
    \cpgeitem{Expliquer pourquoi les fluctuations relatives d'énergie régressent quand la taille du système augmente et associer cette régression au caractère quasi-certain des grandeurs thermodynamiques.}},
  experimentales = {},
  numeriques = {},
}

\cpgeligne{MP/7/2/4}{
  parent = {MP/7/2},
  numero = {7.2.4},
  contenu = {
    \cpgeitem{Système à deux niveaux non dégénérés d’énergies $\pm\epsilon$.}},
  capacites = {
    \cpgeitem{Citer des exemples de systèmes modélisables par un système à deux niveaux.}
    \cpgeitem{Déterminer l’énergie moyenne et la capacité thermique d’un système à deux niveaux.}
    \cpgeitem{Interpréter l'évolution de l'énergie moyenne avec la température, notamment les limites basse et haute température.}
    \cpgeitem{Relier les fluctuations d’énergies à la capacité thermique.}},
  experimentales = {},
  numeriques = {},
}

\cpgeligne{MP/7/2/5}{
  parent = {MP/7/2},
  numero = {7.2.5},
  contenu = {
    \cpgeitem{Énergie moyenne d'équilibre à la température $T$ d’un ensemble de $N$ particules dans un puits de potentiel infini.}},
  capacites = {
    \cpgeitem{Déterminer l’énergie moyenne d’un ensemble de particules à une température donnée, dans la limite où l’énergie de confinement est faible devant l’énergie d’agitation thermique.}
    \cpgeitem{Relier l’expression de l’énergie moyenne en fonction de la température au théorème de l’équipartition de l’énergie.}},
  experimentales = {},
  numeriques = {},
}

\cpgerubrique{MP/7/3}{
  parent = {MP/7},
  numero = {7.3},
  niveau = {2},
  titre = {Capacités thermiques classiques des gaz et des solides},
  chapeau = {},
}

\cpgeligne{MP/7/3/1}{
  parent = {MP/7/3},
  numero = {7.3.1},
  contenu = {
    \cpgeitem{Théorème d'équipartition pour un degré de liberté énergétique indépendant quadratique.}},
  capacites = {
    \cpgeitem{Exploiter la contribution $kT/2$ par degré quadratique à l'énergie moyenne.}},
  experimentales = {},
  numeriques = {},
}

\cpgeligne{MP/7/3/2}{
  parent = {MP/7/3},
  numero = {7.3.2},
  contenu = {
    \cpgeitem{Capacité thermique molaire des gaz classiques dilués monoatomiques et diatomiques.}
    \cpgeitem{Capacité thermique molaire des solides dans le modèle d'Einstein classique : loi de Dulong et Petit.}},
  capacites = {
    \cpgeitem{Dénombrer des degrés de libertés énergétiques quadratiques indépendants et en déduire la capacité thermique molaire d'un système.}},
  experimentales = {},
  numeriques = {},
}

\cpgerubrique{MP/8}{
  parent = {MP},
  numero = {8},
  niveau = {1},
  titre = {Transformations chimiques de la matière : aspects thermodynamiques},
  chapeau = {
    \cpgeitem{Les transformations chimiques de la matière ont été abordées dès le début de la classe de MPSI ; le critère d’évolution spontanée d’un système chimique en transformation y a été présenté sans être démontré. Ce dernier a été remobilisé lors de l’étude des transformations chimiques en solution aqueuse.}
    \cpgeitem{Le but de cette partie est d’une part d’aborder les transferts thermiques et d’autre part d’établir puis exploiter le critère d’évolution spontanée d’un système engagé dans une transformation physico-chimique, ce qui nécessite l’introduction de la fonction enthalpie libre $G$ et du potentiel chimique.}},
}

\cpgerubrique{MP/8/1}{
  parent = {MP/8},
  numero = {8.1},
  niveau = {2},
  titre = {Premier principe de la thermodynamique appliqué aux transformations physico-chimiques},
  chapeau = {
    \cpgeitem{Dans la partie \cpgecettepartie{}, l’étude des transferts thermiques, abordée en première année dans le cadre du cours de physique relatif aux transformation physiques du corps pur, est ici généralisée aux transformations physico-chimiques. Les enthalpies standard de réaction sont considérées comme indépendantes de la température.}
    \cpgeitem{Les notions et contenus sont illustrés à travers des applications liées à la vie quotidienne (contenu calorique des aliments, pouvoirs calorifiques des carburants, etc.), à la recherche (apports des techniques calorimétriques modernes, etc.) ou au domaine industriel.}},
}

\cpgeligne{MP/8/1/1}{
  parent = {MP/8/1},
  numero = {8.1.1},
  contenu = {
    \cpgeitem{État standard. Enthalpie standard de réaction.}
    \cpgeitem{Loi de Hess.}
    \cpgeitem{Enthalpie standard de formation, état standard de référence d'un élément.}},
  capacites = {
    \cpgeitem{Déterminer l'enthalpie standard de réaction à l'aide de tables de données thermodynamiques.}
    \cpgeitem{Associer le signe de l’enthalpie standard de réaction au caractère endothermique ou exothermique de la réaction.}},
  experimentales = {},
  numeriques = {},
}

\cpgeligne{MP/8/1/2}{
  parent = {MP/8/1},
  numero = {8.1.2},
  contenu = {
    \cpgeitem{Effets thermiques pour une transformation monobare : \begin{itemize} \item transfert thermique associé à la transformation chimique en réacteur monobare, isotherme ; \item variation de température en réacteur adiabatique, monobare. \end{itemize}}},
  capacites = {
    \cpgeitem{Prévoir, à partir de données thermodynamiques, le sens et une estimation de la valeur du transfert thermique entre un système, siège d’une transformation physico-chimique et le milieu extérieur.}
    \cpgeitem{Évaluer la température atteinte par un système siège d’une transformation chimique supposée monobare et réalisée dans un réacteur adiabatique.}},
  experimentales = {
    \cpgeitem{Déterminer une enthalpie standard de réaction.}},
  numeriques = {
    \cpgeitem{simuler, à l’aide d’un langage de programmation, l’évolution temporelle de la température pour un système siège d’une transformation adiabatique modélisée par une seule réaction chimique dont les caractéristiques cinétiques et l’enthalpie standard de réaction sont données.}},
}

\cpgerubrique{MP/8/2}{
  parent = {MP/8},
  numero = {8.2},
  niveau = {2},
  titre = {Deuxième principe de la thermodynamique appliqué aux transformations physico-chimiques},
  chapeau = {
    \cpgeitem{Dans la partie \cpgecettepartie{}, on adopte pour les potentiels chimiques une expression générale : $\mu_i = \mu_{i,\text{réf}} +RT \ln(a_i)$ qui fait référence aux activités ai introduites en première année. L'établissement de cette expression est hors programme. L’influence de la pression sur le potentiel chimique d’un constituant en phase condensée pure n’est pas abordée. On se limite aux cas d'une espèce chimique pure, d'une espèce en solution aqueuse très diluée et d'une espèce en mélange de gaz parfaits avec référence à l'état standard.}
    \cpgeitem{Pour le calcul des grandeurs standard de réaction, les enthalpies et entropies standard de réaction sont supposées indépendantes de la température. Les grandeurs standard de réaction permettent la détermination de la valeur, à une température donnée, de la constante thermodynamique d’équilibre $K^{\circ}$ caractéristique d’une réaction, valeur qui était systématiquement fournie en première année. C’est ainsi l’occasion de revenir sur la détermination de la composition d’un système physico-chimique en fin d’évolution.}
    \cpgeitem{La notion d’affinité chimique n’est pas utilisée, le sens d’évolution spontanée d’un système hors d’équilibre, à température et pression fixées, est déterminé par le signe de l’enthalpie libre de réaction $\Delta_{\mathrm{r}} G$. Enfin, l’étude de l’influence de la modification d’un paramètre (pression, température ou composition) sur un système chimique permet d’aborder la problématique de l’optimisation des conditions opératoires d’un procédé chimique.}
    \cpgeitem{Les illustrations et applications sont choisis dans le domaine industriel, dans la vie courante et au niveau du laboratoire.}},
}

\cpgeligne{MP/8/2/1}{
  parent = {MP/8/2},
  numero = {8.2.1},
  contenu = {
    \cpgeitem{Potentiel chimique ; enthalpie libre d’un système chimique. Activité.}},
  capacites = {
    \cpgeitem{Définir le potentiel chimique à l’aide de la fonction enthalpie libre et donner l’expression (admise) du potentiel chimique d’un constituant en fonction de son activité.}
    \cpgeitem{Exprimer l’enthalpie libre d’un système chimique en fonction des potentiels chimiques.}},
  experimentales = {},
  numeriques = {},
}

\cpgeligne{MP/8/2/2}{
  parent = {MP/8/2},
  numero = {8.2.2},
  contenu = {
    \cpgeitem{Enthalpie de réaction, entropie de réaction, enthalpie libre de réaction et grandeurs standard associées.}
    \cpgeitem{Relation entre enthalpie libre de réaction et quotient de réaction ; évolution d’un système chimique.}},
  capacites = {
    \cpgeitem{Justifier qualitativement ou prévoir le signe de l’entropie standard de réaction.}
    \cpgeitem{Relier création d’entropie et enthalpie libre de réaction lors d’une transformation d’un système physico-chimique à pression et température fixées.}
    \cpgeitem{Prévoir le sens d’évolution à pression et température fixées d’un système physico-chimique dans un état donné à l’aide de l’enthalpie libre de réaction.}
    \cpgeitem{Déterminer les grandeurs standard de réaction à partir des tables de données thermodynamiques et de la loi de Hess.}
    \cpgeitem{Déterminer les grandeurs standard de réaction d'une réaction dont l’équation est combinaison linéaire d’autres équations de réaction.}},
  experimentales = {},
  numeriques = {},
}

\cpgeligne{MP/8/2/3}{
  parent = {MP/8/2},
  numero = {8.2.3},
  contenu = {
    \cpgeitem{Constante thermodynamique d’équilibre ; relation de Van ’t Hoff.}},
  capacites = {
    \cpgeitem{Citer et exploiter la relation de Van ’t Hoff. Déterminer la valeur de la constante thermodynamique d’équilibre à une température quelconque.}},
  experimentales = {
    \cpgeitem{Déterminer l’évolution de la valeur d'une constante thermodynamique d'équilibre en fonction de la température.}},
  numeriques = {},
}

\cpgeligne{MP/8/2/4}{
  parent = {MP/8/2},
  numero = {8.2.4},
  contenu = {
    \cpgeitem{État final d’un système : équilibre chimique ou transformation totale.}},
  capacites = {
    \cpgeitem{Déterminer la composition chimique d’un système dans l’état final, en distinguant les cas d’équilibre chimique et de transformation totale, pour une transformation modélisée par une réaction chimique unique.}},
  experimentales = {},
  numeriques = {},
}

\cpgeligne{MP/8/2/5}{
  parent = {MP/8/2},
  numero = {8.2.5},
  contenu = {
    \cpgeitem{Optimisation thermodynamique d’un procédé chimique : \begin{itemize} \item par modification de la valeur de $K^{\circ}$ ; \item par modification de la valeur du quotient de réaction. \end{itemize}}},
  capacites = {
    \cpgeitem{Identifier les paramètres d’influence et leur contrôle pour optimiser une synthèse ou minimiser la formation d’un produit secondaire indésirable.}},
  experimentales = {},
  numeriques = {},
}

\cpgerubrique{MP/9}{
  parent = {MP},
  numero = {9},
  niveau = {1},
  titre = {Aspects thermodynamique et cinétique de l’électrochimie},
  chapeau = {
    \cpgeitem{Les aspects thermodynamiques et cinétiques des réactions d’oxydo-réduction sont appliqués notamment à la corrosion d’une part et aux dispositifs électrochimiques que sont les piles et les accumulateurs d’autre part. L’illustration des notions gagne à s’appuyer sur des applications concrètes comme par exemple la mise en œuvre de capteurs électrochimiques dans l’analyse de l’eau, de l’air ou d’effluents.}
    \cpgeitem{L’approche de l’électrochimie proposée ici privilégie les raisonnements qualitatifs et les aspects expérimentaux, plutôt que les développements théoriques et formels.}},
}

\cpgerubrique{MP/9/1}{
  parent = {MP/9},
  numero = {9.1},
  niveau = {2},
  titre = {Étude thermodynamique des réactions d’oxydo-réduction},
  chapeau = {
    \cpgeitem{La partie \cpgecettepartie{} se fonde sur les acquis de première année relatifs à l’étude des réactions d’oxydo-réduction et des piles, ainsi que sur la partie de thermodynamique chimique de seconde année pour relier les grandeurs thermodynamiques aux potentiels.}},
}

\cpgeligne{MP/9/1/1}{
  parent = {MP/9/1},
  numero = {9.1.1},
  contenu = {
    \cpgeitem{Relation entre enthalpie libre de réaction et potentiels des couples mis en jeu dans une réaction d’oxydo-réduction.}},
  capacites = {
    \cpgeitem{Citer et exploiter la relation entre l’enthalpie libre de réaction et les potentiels des couples mis en jeu dans une réaction d’oxydo-réduction.}},
  experimentales = {},
  numeriques = {},
}

\cpgeligne{MP/9/1/2}{
  parent = {MP/9/1},
  numero = {9.1.2},
  contenu = {
    \cpgeitem{Relation entre enthalpie libre standard de réaction et potentiels standard des couples impliqués.}},
  capacites = {
    \cpgeitem{Déterminer l’enthalpie libre standard d’une réaction d’oxydo-réduction à partir des potentiels standard des couples.}
    \cpgeitem{Déterminer la valeur du potentiel standard d’un couple d’oxydo-réduction à partir de données thermodynamiques.}},
  experimentales = {},
  numeriques = {},
}

\cpgerubrique{MP/9/2}{
  parent = {MP/9},
  numero = {9.2},
  niveau = {2},
  titre = {Étude cinétique des réactions d’oxydo-réduction : courbe courant-potentiel},
  chapeau = {
    \cpgeitem{La partie \cpgecettepartie{} se fonde sur les acquis de cinétique chimique de première année et les prolongent par le tracé et l’exploitation de courbes courant-potentiel.}
    \cpgeitem{Les courbes courant-potentiel, dont le tracé est proposé en capacité expérimentale, sont un outil essentiel dans la compréhension et la modélisation des systèmes électrochimiques.}
    \cpgeitem{L’écart entre le potentiel d’une électrode et son potentiel d’équilibre est appelé surpotentiel plutôt que surtension pour des raisons pédagogiques, en cohérence avec le vocabulaire anglo-saxon correspondant.}},
}

\cpgeligne{MP/9/2/1}{
  parent = {MP/9/2},
  numero = {9.2.1},
  contenu = {
    \cpgeitem{Courbes courant-potentiel sur une électrode en régime stationnaire : \begin{itemize} \item surpotentiel ; \item systèmes rapides et systèmes lents ; \item nature de l’électrode ; \item courant de diffusion limite ; \item vagues successives ; \item domaine d’inertie électrochimique du solvant. \end{itemize}}},
  capacites = {
    \cpgeitem{Décrire le montage à trois électrodes permettant de tracer des courbes courant-potentiel.}
    \cpgeitem{Relier vitesse de réaction électrochimique et intensité du courant.}
    \cpgeitem{Identifier le caractère lent ou rapide d’un système à partir des courbes courant-potentiel. Identifier les espèces électroactives pouvant donner lieu à une limitation en courant par diffusion.}
    \cpgeitem{Relier qualitativement ou quantitativement, à partir de relevés expérimentaux, l'intensité du courant de diffusion limite à la concentration du réactif et à la surface immergée de l'électrode.}
    \cpgeitem{Tracer l’allure de courbes courant-potentiel de branches d’oxydation ou de réduction à partir de données fournies, de potentiels standard, concentrations et surpotentiels.}},
  experimentales = {
    \cpgeitem{Tracer et exploiter des courbes courant-potentiel.}},
  numeriques = {},
}

\cpgerubrique{MP/9/3}{
  parent = {MP/9},
  numero = {9.3},
  niveau = {2},
  titre = {Stockage et conversion d’énergie chimique dans des dispositifs électrochimiques},
  chapeau = {
    \cpgeitem{La partie \cpgecettepartie{} s’appuie sur les courbes courant-potentiel pour étudier le fonctionnement des piles et leur recharge, ainsi que les électrolyseurs. Ces courbes permettent en effet de déterminer différentes caractéristiques : réactions aux électrodes, tension à vide, tension à imposer pour une recharge, etc.}},
}

\cpgeligne{MP/9/3/1}{
  parent = {MP/9/3},
  numero = {9.3.1},
  contenu = {
    \cpgeitem{\cpgenotion{Conversion d’énergie chimique en énergie électrique : fonctionnement des piles.}}
    \cpgeitem{Transformations spontanées et réaction modélisant le fonctionnement d’une pile électrochimique.}},
  capacites = {
    \cpgeitem{Établir l’inégalité reliant la variation d’enthalpie libre et le travail électrique.}
    \cpgeitem{Relier la tension à vide d’une pile et l’enthalpie libre de la réaction modélisant son fonctionnement.}
    \cpgeitem{Déterminer la capacité électrique d’une pile.}},
  experimentales = {},
  numeriques = {},
}

\cpgeligne{MP/9/3/2}{
  parent = {MP/9/3},
  numero = {9.3.2},
  contenu = {
    \cpgeitem{Courbes courant-potentiel et fonctionnement d’une pile électrochimique.}},
  capacites = {
    \cpgeitem{Exploiter les courbes courant-potentiel pour rendre compte du fonctionnement d’une pile électrochimique et tracer sa caractéristique.}
    \cpgeitem{Citer les paramètres influençant la résistance interne d’une pile électrochimique.}},
  experimentales = {},
  numeriques = {},
}

\cpgeligne{MP/9/3/3}{
  parent = {MP/9/3},
  numero = {9.3.3},
  contenu = {
    \cpgeitem{\cpgenotion{Conversion d’énergie électrique en énergie chimique.}}
    \cpgeitem{Transformations forcées lors d’une électrolyse et de la recharge d’un accumulateur.}},
  capacites = {
    \cpgeitem{Exploiter les courbes courant-potentiel pour rendre compte du fonctionnement d’un électrolyseur et prévoir la valeur de la tension minimale à imposer.}
    \cpgeitem{Exploiter les courbes courant-potentiel pour justifier les contraintes (purification de la solution électrolytique, choix des électrodes) dans la recharge d’un accumulateur.}
    \cpgeitem{Déterminer la masse de produit formé pour une durée et des conditions données d’électrolyse.}
    \cpgeitem{Déterminer un rendement faradique à partir d’informations fournies concernant le dispositif étudié.}},
  experimentales = {},
  numeriques = {},
}

\cpgeligne{MP/9/3/4}{
  parent = {MP/9/3},
  numero = {9.3.4},
  contenu = {
    \cpgeitem{Stockage et conversion d’énergie chimique.}},
  capacites = {},
  experimentales = {
    \cpgeitem{Étudier le fonctionnement d’une pile ou d’un électrolyseur pour effectuer des bilans de matière et des bilans électriques.}},
  numeriques = {},
}

\cpgerubrique{MP/9/4}{
  parent = {MP/9},
  numero = {9.4},
  niveau = {2},
  titre = {Corrosion humide et électrochimique},
  chapeau = {
    \cpgeitem{La lutte contre la corrosion est un enjeu économique et la compréhension des phénomènes de corrosion et des facteurs influençant cette corrosion est essentielle pour effectuer des choix de méthodes de protection. La partie «\cpgecettepartie{} exploite les courbes courant-potentiel pour interpréter les phénomènes de corrosion, de protection et de passivation. On se limite à la corrosion uniforme et à la corrosion galvanique de deux métaux en contact. Les tracés de diagrammes de Tafel ou d’Evans sont hors-programme.}},
}

\cpgeligne{MP/9/4/1}{
  parent = {MP/9/4},
  numero = {9.4.1},
  contenu = {
    \cpgeitem{Corrosion uniforme en milieu acide ou en milieu neutre oxygéné : potentiel de corrosion, courant de corrosion.}
    \cpgeitem{Corrosion d’un système de deux métaux en contact.}},
  capacites = {
    \cpgeitem{Positionner un potentiel de corrosion sur un tracé de courbes courant-potentiel.}
    \cpgeitem{Interpréter le phénomène de corrosion uniforme d’un métal ou de deux métaux en contact en utilisant des courbes courant-potentiel ou d’autres données expérimentales, thermodynamiques et cinétiques.}
    \cpgeitem{Citer des facteurs favorisant la corrosion.}},
  experimentales = {},
  numeriques = {},
}

\cpgeligne{MP/9/4/2}{
  parent = {MP/9/4},
  numero = {9.4.2},
  contenu = {
    \cpgeitem{Protection contre la corrosion : \begin{itemize} \item revêtement ; \item anode sacrificielle ; \item protection électrochimique par courant imposé. \end{itemize}}},
  capacites = {
    \cpgeitem{Exploiter des tracés de courbes courant-potentiel pour expliquer qualitativement : \begin{itemize} \item la qualité de la protection par un revêtement métallique ; \item le fonctionnement d’une anode sacrificielle. \end{itemize}}},
  experimentales = {},
  numeriques = {},
}

\cpgeligne{MP/9/4/3}{
  parent = {MP/9/4},
  numero = {9.4.3},
  contenu = {
    \cpgeitem{Passivation.}},
  capacites = {
    \cpgeitem{Interpréter le phénomène de passivation sur une courbe courant-potentiel.}},
  experimentales = {
    \cpgeitem{Mettre en évidence le phénomène de corrosion et de protection et des facteurs l’influençant.}},
  numeriques = {},
}

\cpgerubrique{MPA/A1}{
  parent = {MP},
  numero = {A1},
  niveau = {1},
  titre = {Matériel},
  chapeau = {
    \cpgeitem{Cette liste complète celle donnée en annexe 1 du programme de physique-chimie de la classe de MPSI. À elles deux, ces listes regroupent le matériel que les étudiants doivent savoir utiliser avec l’aide d’une notice simplifiée. Une utilisation de matériel hors de ces listes lors d’épreuves d’évaluation n’est pas exclue, mais elle doit obligatoirement s’accompagner d’une aide.}},
}

\cpgeligne{MPA/A1/1}{
  parent = {MPA/A1},
  numero = {A1.1},
  contenu = {
    \cpgeitem{Domaine optique}},
  capacites = {
    \cpgeitem{\begin{itemize} \item Polariseur. \item Interféromètre de Michelson motorisé. \item Capteur photographique numérique. \item Spectromètre à fibre optique. \end{itemize}}},
  experimentales = {},
  numeriques = {},
}

\cpgeligne{MPA/A1/2}{
  parent = {MPA/A1},
  numero = {A1.2},
  contenu = {
    \cpgeitem{Domaine électrique}},
  capacites = {
    \cpgeitem{\begin{itemize} \item Oscilloscope numérique avec analyseur de spectre. \item Carte d'acquisition dont l'API est publiée. \item Microcontrôleur. \item Émetteur et récepteur dans le domaine des ondes centimétriques. \end{itemize}}},
  experimentales = {},
  numeriques = {},
}

\cpgeligne{MPA/A1/3}{
  parent = {MPA/A1},
  numero = {A1.3},
  contenu = {
    \cpgeitem{Domaine de la chimie}},
  capacites = {
    \cpgeitem{\begin{itemize} \item Calorimètre. \item Électrode de référence. \item Électrolyseur et électrodes. \end{itemize}}},
  experimentales = {},
  numeriques = {},
}

\cpgerubrique{MPA/A2}{
  parent = {MP},
  numero = {A2},
  niveau = {1},
  titre = {Outils mathématiques},
  chapeau = {
    \cpgeitem{Les outils mathématiques dont la maîtrise est nécessaire à la mise en œuvre du programme de physique de la classe de MP sont d’une part ceux qui figurent dans l’annexe 2 du programme de la classe de MPSI et d’autre part ceux qui figurent dans la liste ci-dessous.}},
}

\cpgerubrique{MPA/A2/1}{
  parent = {MPA/A2},
  numero = {A2.1},
  niveau = {2},
  titre = {Analyse vectorielle},
  chapeau = {
    \cpgeitem{Le thème \cpgecettepartie{} prolonge l’outil « gradient » abordé en première année, en introduisant de nouveaux opérateurs : seules les expressions des opérateurs en coordonnées cartésiennes sont exigibles. Les expressions des opérateurs en coordonnées cylindriques et sphériques et les formules d’analyse vectorielle ne sont pas exigibles ; elles doivent donc être systématiquement rappelées.}},
}

\cpgeligne{MPA/A2/1/1}{
  parent = {MPA/A2/1},
  numero = {A2.1.1},
  contenu = {
    \cpgeitem{Gradient.}},
  capacites = {
    \cpgeitem{Citer le lien entre le gradient et la différentielle.}
    \cpgeitem{Exprimer les composantes du gradient en coordonnées cartésiennes.}
    \cpgeitem{Utiliser le fait que le gradient d’une fonction $f$ est perpendiculaire aux surfaces iso-$f$ et orienté dans le sens des valeurs de $f$ croissantes.}},
  experimentales = {},
  numeriques = {},
}

\cpgeligne{MPA/A2/1/2}{
  parent = {MPA/A2/1},
  numero = {A2.1.2},
  contenu = {
    \cpgeitem{Divergence.}},
  capacites = {
    \cpgeitem{Citer et utiliser le théorème d’Ostrogradski. Exprimer la divergence en coordonnées cartésiennes.}},
  experimentales = {},
  numeriques = {},
}

\cpgeligne{MPA/A2/1/3}{
  parent = {MPA/A2/1},
  numero = {A2.1.3},
  contenu = {
    \cpgeitem{Rotationnel.}},
  capacites = {
    \cpgeitem{Citer et utiliser le théorème de Stokes. Exprimer le rotationnel en coordonnées cartésiennes.}},
  experimentales = {},
  numeriques = {},
}

\cpgeligne{MPA/A2/1/4}{
  parent = {MPA/A2/1},
  numero = {A2.1.4},
  contenu = {
    \cpgeitem{Laplacien d'un champ scalaire.}},
  capacites = {
    \cpgeitem{Définir le laplacien à l’aide de la divergence et du gradient.}
    \cpgeitem{Exprimer le laplacien en coordonnées cartésiennes.}},
  experimentales = {},
  numeriques = {},
}

\cpgeligne{MPA/A2/1/5}{
  parent = {MPA/A2/1},
  numero = {A2.1.5},
  contenu = {
    \cpgeitem{Laplacien d'un champ de vecteurs.}},
  capacites = {
    \cpgeitem{Exprimer le laplacien d’un champ de vecteurs en coordonnées cartésiennes.}
    \cpgeitem{Utiliser la formule d’analyse vectorielle : $\cpgerot(\cpgerot(\cpgevec A)) = \cpgegrad(\cpgediv\cpgevec A) - \Delta\cpgevec A$.}},
  experimentales = {},
  numeriques = {},
}

\cpgeligne{MPA/A2/1/6}{
  parent = {MPA/A2/1},
  numero = {A2.1.6},
  contenu = {
    \cpgeitem{Cas des champs proportionnels à $\exp(\cpgeii\omega t- \cpgeii\cpgevec k\cdot\cpgevec r)$ ou à $\exp(\cpgeii\cpgevec k\cdot\cpgevec r-\cpgeii\omega t)$.}},
  capacites = {
    \cpgeitem{Exprimer l'action des opérateurs d'analyse vectorielle sur un tel champ à l'aide du vecteur $\cpgeii\cpgevec k$.}},
  experimentales = {},
  numeriques = {},
}

\cpgerubrique{MPA/A2/2}{
  parent = {MPA/A2},
  numero = {A2.2},
  niveau = {2},
  titre = {Analyse de Fourier},
  chapeau = {
    \cpgeitem{Le thème \cpgecettepartie{} prolonge l’étude de l’outil « séries de Fourier » abordée en MPSI et réutilisée en classe de MP. On étend la décomposition d’un signal périodique comme somme de ses harmoniques à l’expression d’un signal non périodique sous forme d’une intégrale (synthèse spectrale) ; aucun résultat n’est exigible. On souligne en revanche la relation liant en ordre de grandeur la largeur spectrale $\Delta f$ et la durée caractéristique $\Delta t$ d’un signal non périodique.}},
}

\cpgeligne{MPA/A2/2/1}{
  parent = {MPA/A2/2},
  numero = {A2.2.1},
  contenu = {
    \cpgeitem{Décomposition d’une fonction périodique en série de Fourier.}},
  capacites = {
    \cpgeitem{Utiliser un développement en série de Fourier fourni.}
    \cpgeitem{Utiliser un raisonnement par superposition.}},
  experimentales = {},
  numeriques = {},
}

\cpgeligne{MPA/A2/2/2}{
  parent = {MPA/A2/2},
  numero = {A2.2.2},
  contenu = {
    \cpgeitem{Synthèse spectrale d’un signal non périodique.}},
  capacites = {
    \cpgeitem{Utiliser un raisonnement par superposition.}
    \cpgeitem{Citer et utiliser la relation liant en ordre de grandeur la largeur spectrale $\Delta f$ et la durée caractéristique $\Delta t$ d’un signal non périodique.}},
  experimentales = {},
  numeriques = {},
}

\cpgerubrique{MPA/A2/3}{
  parent = {MPA/A2},
  numero = {A2.3},
  niveau = {2},
  titre = {Équations aux dérivées partielles},
  chapeau = {
    \cpgeitem{Dans le thème \cpgecettepartie{}, aucune méthode générale d’étude n’est exigible : on se limite à chercher des solutions d’une forme donnée par substitution, menant ainsi soit à des équations différentielles classiques, soit à une relation de dispersion. L’accent est mis sur le rôle des conditions aux limites.}},
}

\cpgeligne{MPA/A2/3/1}{
  parent = {MPA/A2/3},
  numero = {A2.3.1},
  contenu = {
    \cpgeitem{Exemples d’équations aux dérivées partielles : équation de Laplace, équation de diffusion, équation de d’Alembert, équation de Schrödinger.}},
  capacites = {
    \cpgeitem{Identifier une équation aux dérivées partielles connue.}
    \cpgeitem{Transposer une solution familière dans un domaine de la physique à un autre domaine.}
    \cpgeitem{Obtenir des solutions de forme donnée par substitution.}
    \cpgeitem{Utiliser des conditions initiales et des conditions aux limites.}},
  experimentales = {},
  numeriques = {},
}

\cpgerubrique{MPA/A2/4}{
  parent = {MPA/A2},
  numero = {A2.4},
  niveau = {2},
  titre = {Calcul différentiel},
  chapeau = {
    \cpgeitem{Les capacités relatives à la notion de différentielle d’une fonction de plusieurs variables sont limitées à l’essentiel, elles sont mobilisées principalement dans le cours de thermodynamique.}},
}

\cpgeligne{MPA/A2/4/1}{
  parent = {MPA/A2/4},
  numero = {A2.4.1},
  contenu = {
    \cpgeitem{Différentielle d’une fonction de plusieurs variables. Dérivée partielle. Théorème de Schwarz.}},
  capacites = {
    \cpgeitem{Connaître l’expression de la différentielle en fonction des dérivées partielles.}
    \cpgeitem{Identifier la valeur d’une dérivée partielle, l’expression de la différentielle étant donnée.}
    \cpgeitem{Utiliser le théorème de Schwarz (admis).}},
  experimentales = {},
  numeriques = {},
}

\cpgerubrique{MPA/A3}{
  parent = {MP},
  numero = {A3},
  niveau = {1},
  titre = {Outils numériques},
  chapeau = {
    \cpgeitem{La prise en compte de capacités de codage en langage Python dans la formation des étudiants inclut l’utilisation de fonctions extraites de diverses bibliothèques. Elle vise à une meilleure appréhension des principes mis en œuvre par les différents logiciels de traitement des données dont l’utilisation est par ailleurs toujours recommandée. Elle a aussi pour objectif de mobiliser ces capacités dans un contexte concret, celui de la physique. Cette formation par le codage permet également de développer des capacités utiles à la physique comme le raisonnement, la logique ou la décomposition d’un problème complexe en étapes plus simples.}},
}

\cpgeligne{MPA/A3/1}{
  parent = {MPA/A3},
  numero = {A3.1.1},
  contenu = {
    \cpgeitem{Équation de diffusion à une dimension.}},
  capacites = {
    \cpgeitem{Mettre en œuvre une méthode des différences finies explicite pour résoudre l’équation de diffusion à une dimension en régime variable.}},
  experimentales = {},
  numeriques = {},
}

